% From NumPy and SciPy to simplydata: written by hand, read by docs/make_manual.py (% From NumPy and SciPy to simplydata: written by hand, read by docs/make_manual.py (% From NumPy and SciPy to simplydata: written by hand, read by docs/make_manual.py (% From NumPy and SciPy to simplydata: written by hand, read by docs/make_manual.py (\input{from_numpy}).
% The topics follow the order of the NumPy absolute beginners guide and the NumPy user guide (fundamentals), then the
% SciPy user guide, used as a checklist only: the text and the examples are our own, and each section links to the
% official page. Every example in a py box is run by tests/test_from_numpy_chapter.py (with "import simplydata as
% sdata" in front), so it cannot go out of date.

\unchapter{From NumPy and SciPy to simplydata}{from-numpy}

This chapter is for you if you already know NumPy or SciPy, or if you want to learn them after
simplydata. It follows the topics of the official guides, in their order: the
\href{https://numpy.org/doc/stable/user/absolute_beginners.html}{NumPy absolute beginners guide} and
the \href{https://numpy.org/doc/stable/user/basics.html}{NumPy fundamentals}, then the
\href{https://docs.scipy.org/doc/scipy/tutorial/index.html}{SciPy user guide}. For each topic it says
what the job is, which simplydata words do it, and a short example. Some jobs are left to NumPy and
SciPy on purpose. For those, the chapter shows how to hand the numbers over and how to bring the
answer back.

Every example starts with \code{import simplydata as sdata}. It is written once here, not in each box.

\unsection{Five things to know first}{from-numpy:first}

\begin{enumerate}
  \item \textbf{Three makers, one family.} \code{sdata.number\_list} makes a list of numbers (NumPy's
    1-D array), \code{sdata.number\_grid} makes rows and columns (a 2-D array) and
    \code{sdata.number\_block} makes numbers in any number of directions, each with a name (an n-D
    array, or an xarray DataArray). The words of numbers work on grids and blocks too.
  \item \textbf{Positions count from 1.} \code{numbers.get\_number\_at(1)} is the first number. NumPy
    counts from 0, so the sentence a word prints gives NumPy's position too.
  \item \textbf{A word changes the object, and \code{undo} takes it back.} \code{numbers.take\_log()}
    changes \code{numbers} itself and prints what it did. \code{numbers.undo()} brings back the numbers
    from before. A word that finds something (an average, a determinant) gives it back and changes
    nothing.
  \item \textbf{Empty values stay empty.} NumPy writes a missing value as \code{nan}. simplydata calls it
    an empty value, says how many there are, and leaves them out of averages and totals, as NumPy's
    \code{nanmean} does.
  \item \textbf{The door is always open.} \code{array = numbers.as\_numpy()} gives a plain NumPy array,
    for any NumPy or SciPy call. \code{sdata.number\_list(array)}, \code{sdata.number\_grid(array)} and
    \code{sdata.number\_block(array)} bring the answer back.
\end{enumerate}

\begin{lstlisting}[style=py]
numbers = sdata.number_list([2.5, 3.1, 4.8, 5.0, 6.2])
array = numbers.as_numpy()               # a NumPy array: use any NumPy or SciPy code here
import numpy as np
steps = np.diff(array)
back = sdata.number_list(steps)          # and back to simplydata
back.find_average()
\end{lstlisting}

Do you know a NumPy, SciPy or SymPy name, but not the simplydata word? Ask:
\code{sdata.help("np.mean")}, \code{sdata.help("np.linalg.solve")}, \code{sdata.help("rng.normal")}
or \code{sdata.help("scipy.stats.gamma")}. The name can have \code{np.}, \code{np.linalg.},
\code{scipy.}, \code{stats.} or \code{optimize.} in front, and the answer is the word for that package:
\code{np.mean} is the average of numbers, but a bare \code{mean} is the average of a table column.

\unsection{Making numbers, grids and blocks}{from-numpy:making}

\textbf{The job.} Make an array: from a list, as a run of evenly spaced numbers, filled with one
value, or of a given shape. In the guides:
\href{https://numpy.org/doc/stable/user/absolute_beginners.html}{the absolute beginners guide}, ``How
to create a basic array'', and \href{https://numpy.org/doc/stable/user/basics.creation.html}{Array
creation}.

\textbf{In simplydata.} \code{sdata.number\_list([...])} is \code{np.array}; \code{sdata.number\_list(0,
to=1, how\_many=11)} is \code{np.linspace} and \code{step=} is \code{np.arange}.
\code{sdata.number\_grid(rows=2, columns=3, filled\_with=0)} is \code{np.zeros} (and \code{np.full});
\code{sdata.number\_block(of\_shape=(2, 3, 4), filled\_with=1)} is \code{np.ones} with any shape.
Special grids come from \code{kind=}: \code{sdata.number\_grid(kind="identity", size=3)} is \code{np.eye},
\code{kind="diagonal"} is \code{np.diag}, and \code{"toeplitz"}, \code{"hilbert"} or \code{"pascal"} are the
grids of \code{scipy.linalg}. \code{sdata.make\_coordinate\_grids(x=..., y=...)} is \code{np.meshgrid}.

\begin{lstlisting}[style=py]
temperatures = sdata.number_list([8.3, 10.3, 11.0, 9.2, 7.5])
steps = sdata.number_list(0, to=1, how_many=11)          # 0, 0.1, ..., 1
every_five = sdata.number_list(0, to=50, step=5)
grid = sdata.number_grid([[1, 2, 3], [4, 5, 6]])
zeros = sdata.number_grid(rows=2, columns=3, filled_with=0)
block = sdata.number_block(of_shape=(2, 3, 4), filled_with=1)
print(grid)
grid = sdata.number_grid(kind="identity", size=3)       # NumPy's eye
a, b = sdata.make_coordinate_grids(x=[0, 1, 2], y=[10, 20])    # NumPy's meshgrid
\end{lstlisting}

\unsection{Shape and size}{from-numpy:shape}

\textbf{The job.} Know how many numbers there are, how many rows and columns, and how many directions
(NumPy's \code{size}, \code{shape} and \code{ndim}). In the guide:
\href{https://numpy.org/doc/stable/user/absolute_beginners.html}{the absolute beginners guide}, ``How
do you know the shape and size of an array?''.

\textbf{In simplydata.} \code{numbers.count()} is \code{size}; \code{grid.find\_shape()} and
\code{block.find\_shape()} are \code{shape}; \code{block.count\_directions()} is \code{ndim}.
\code{grid.count\_rows()} and \code{grid.count\_columns()} say one size each.

\begin{lstlisting}[style=py]
grid = sdata.number_grid([[1, 2, 3], [4, 5, 6]])
grid.find_shape()                        # (2, 3)
grid.count_rows()
block = sdata.number_block(of_shape=(2, 3, 4), filled_with=0)
block.count_directions()                 # 3
\end{lstlisting}

\unsection{Picking numbers out}{from-numpy:indexing}

\textbf{The job.} Take one number, a slice, a row or a column, or the numbers that pass a test
(NumPy's indexing and slicing, and boolean indexing). In the guides:
\href{https://numpy.org/doc/stable/user/absolute_beginners.html}{the absolute beginners guide},
``Indexing and slicing'', and \href{https://numpy.org/doc/stable/user/basics.indexing.html}{Indexing
on ndarrays}.

\textbf{In simplydata.} \code{get\_number\_at} is \code{a[i]} (from 1); \code{keep\_numbers\_by\_position(2,
to=4)} is a slice and keeps both ends; \code{grid.get\_row} and \code{grid.get\_column} are \code{a[i, :]}
and \code{a[:, j]}. \code{keep\_each\_number\_that(is\_above=5)} is \code{a[a > 5]}.
\code{grid.get\_part("diagonal")} and \code{grid.get\_part("upper triangle")} are \code{np.diagonal} and
\code{np.triu}; \code{find\_positions\_where(is\_above=5)} is \code{np.nonzero} (counting from 1).

\begin{lstlisting}[style=py]
numbers = sdata.number_list([4, 8, 15, 16, 23, 42])
numbers.get_number_at(1)                 # 4, the first number (NumPy's a[0])
numbers.keep_numbers_by_position(2, to=4)
print(numbers)                           # 8, 15, 16
numbers.undo()
numbers.keep_each_number_that(is_above=10)
grid = sdata.number_grid([[1, 2, 3], [4, 5, 6]])
grid.get_row(2)
grid.get_column(3)
grid.get_part("diagonal")                # 1 and 5
\end{lstlisting}

\unsection{Maths on every number}{from-numpy:maths}

\textbf{The job.} Add, multiply, take logs and square roots of every number at once (NumPy's universal
functions), also between two arrays of the same length. In the guides:
\href{https://numpy.org/doc/stable/user/absolute_beginners.html}{the absolute beginners guide}, ``Basic
array operations'', and \href{https://numpy.org/doc/stable/user/basics.ufuncs.html}{Universal
functions}.

\textbf{In simplydata.} \code{add}, \code{subtract}, \code{times}, \code{divide\_by},
\code{raise\_to\_power}, \code{take\_square\_root}, \code{take\_log}, \code{take\_exp},
\code{take\_sin}, \code{round\_to} and their family. Each takes one number or numbers of the same
length. \code{+}, \code{-}, \code{*} and \code{/} work too and give new numbers.
Every other function of NumPy and \code{scipy.special} is one family word:
\code{numbers.take\_function("hyperbolic sine")} is \code{np.sinh}, and
\code{take\_function("bessel J", order=1)} is \code{special.jv}. The library names work too, and
\code{numbers.take\_function()} with no name lists them all. \code{work\_on\_bits} does the bit operations
(\code{np.bitwise\_and} and the rest). \code{turn\_into\_changes} is \code{np.diff},
\code{turn\_into\_slopes} is \code{np.gradient} and \code{find\_product} is \code{np.prod}.

\begin{lstlisting}[style=py]
prices = sdata.number_list([1.20, 3.50, 2.25])
prices.times(1.25)                       # every price with 25 % added
prices.round_to(with_decimals=2)
counts = sdata.number_list([10, 100, 1000])
counts.take_log(base=10)
both = sdata.number_list([1, 2, 3]) + sdata.number_list([10, 20, 30])
print(both)
numbers = sdata.number_list([0.5, 1.0, 2.0])
numbers.take_function("hyperbolic sine")                 # NumPy's sinh
numbers.undo()
numbers.take_function("bessel J", order=1)               # SciPy's special.jv
numbers = sdata.number_list([3, 5, 9, 10])
numbers.turn_into_changes()                              # 2, 4, 1: NumPy's diff
\end{lstlisting}

\unsection{Stretching to the same shape}{from-numpy:broadcasting}

\textbf{The job.} Work on two arrays of different shapes by stretching the smaller one (NumPy's
broadcasting). In the guide: \href{https://numpy.org/doc/stable/user/basics.broadcasting.html}%
{Broadcasting}.

\textbf{In simplydata.} The maths words stretch by themselves, as NumPy does, and print a note when
they do. \code{block.show\_how\_it\_lines\_up\_with(other)} explains how two blocks line up before you
use them, and \code{sdata.stretch\_to\_same\_shape(a, b)} gives both at the same shape.
\code{sdata.hide\_stretching\_notes()} turns the notes off.

\begin{lstlisting}[style=py]
a = sdata.number_block(of_shape=(3, 4), filled_with=1)
b = sdata.number_block(of_shape=(1, 4), filled_with=2)
a.show_how_it_lines_up_with(b)
a_big, b_big = sdata.stretch_to_same_shape(a, b)
b_big.find_shape()
\end{lstlisting}

\unsection{Summaries}{from-numpy:summaries}

\textbf{The job.} The total, average, smallest and biggest number, the standard deviation and
percentiles, of all numbers or along one direction (NumPy's aggregation functions with
\code{axis=}). In the guide: \href{https://numpy.org/doc/stable/user/absolute_beginners.html}{the
absolute beginners guide}, ``More useful array operations''.

\textbf{In simplydata.} \code{find\_total}, \code{find\_average}, \code{find\_median},
\code{find\_smallest}, \code{find\_biggest}, \code{find\_standard\_deviation} and
\code{find\_percentile}. On a block, \code{find\_average(over="month")} averages along a named
direction, as \code{axis=} does with a number. The standard deviation is the sample one (NumPy's
\code{ddof=1}), as in statistics books.

\begin{lstlisting}[style=py]
numbers = sdata.number_list([8.3, 10.3, 11.0, 9.2, 7.5])
numbers.find_average()
numbers.find_standard_deviation()        # sample standard deviation (NumPy's ddof=1)
numbers.find_percentile(90)
block = sdata.number_block(of_shape=(2, 3), filled_with=1, directions_named=["year", "month"])
block.find_average(over="month")
\end{lstlisting}

\unsection{Changing the shape}{from-numpy:reshaping}

\textbf{The job.} Reshape, turn rows into columns, join arrays and split them, add or remove a
direction (NumPy's \code{reshape}, \code{transpose}, \code{concatenate}, \code{split},
\code{expand\_dims} and \code{squeeze}). In the guide:
\href{https://numpy.org/doc/stable/user/absolute_beginners.html}{the absolute beginners guide}, ``Can
you reshape an array?'' and ``Transposing and reshaping a matrix''.

\textbf{In simplydata.} \code{numbers.make\_number\_grid(rows=, columns=)} makes a grid from numbers;
\code{block.change\_shape\_to(...)} changes a block's own shape. \code{grid.swap\_rows\_and\_columns()}
is \code{.T}. \code{block.add\_numbers\_from(other, over=)} joins, \code{block.split\_into\_pieces}
splits, and \code{add\_direction} and \code{remove\_directions\_of\_size\_one} are \code{expand\_dims}
and \code{squeeze}.

\begin{lstlisting}[style=py]
numbers = sdata.number_list([1, 2, 3, 4, 5, 6])
grid = numbers.make_number_grid(rows=2, columns=3)
grid.swap_rows_and_columns()
grid.find_shape()                        # (3, 2)
block = sdata.number_block(of_shape=(2, 6), filled_with=0, directions_named=["row", "column"])
block.change_shape_to(2, 3, 2)
block.find_shape()
\end{lstlisting}

\unsection{Sorting, unique numbers and counts}{from-numpy:sorting}

\textbf{The job.} Sort, keep each different number once, and count how often each comes (NumPy's
\code{sort}, \code{unique} and \code{unique\_counts}). In the guide:
\href{https://numpy.org/doc/stable/user/absolute_beginners.html}{the absolute beginners guide},
``Adding, removing, and sorting elements'' and ``How to get unique items and counts''.

\textbf{In simplydata.} \code{sort} (\code{biggest\_first=True} for the other way),
\code{remove\_duplicate\_numbers} and \code{count\_each}, which gives a table of each number and its
count. \code{sort(show="positions")} is \code{np.argsort}, and \code{count\_in\_bins(how\_many\_bins=3)} is
\code{np.histogram} as numbers.

\begin{lstlisting}[style=py]
numbers = sdata.number_list([3, 1, 2, 3, 3, 1])
numbers.count_each()
numbers.sort(biggest_first=True)
numbers.remove_duplicate_numbers()
print(numbers)
numbers = sdata.number_list([3, 1, 2, 3, 3, 1])
numbers.count_in_bins(how_many_bins=3)
\end{lstlisting}

\unsection{Yes or no, and conditions}{from-numpy:conditions}

\textbf{The job.} Ask a question of every number and use the answers: count them, keep the numbers
where the answer is yes, or replace them (NumPy's comparisons, boolean masks, \code{where},
\code{count\_nonzero}, \code{any} and \code{all}). In the guide:
\href{https://numpy.org/doc/stable/user/basics.indexing.html}{Indexing on ndarrays}, ``Boolean array
indexing''.

\textbf{In simplydata.} \code{turn\_into\_yes\_or\_no(is\_above=)} is the mask;
\code{keep\_numbers\_where}, \code{replace\_numbers\_where} and \code{count\_numbers\_that} use it, and
\code{check\_if\_any\_number} and \code{check\_if\_all\_numbers} are \code{any} and \code{all}.

\begin{lstlisting}[style=py]
rain = sdata.number_list([0.0, 2.5, 0.0, 12.1, 4.4, 0.0])
rain.count_numbers_that(is_above=0)      # 3 days with rain
wet = rain.turn_into_yes_or_no(is_above=0)
rain.check_if_any_number(is_above=10)
rain.replace_numbers_where(wet, with_value=1)
print(rain)
\end{lstlisting}

\unsection{Empty values}{from-numpy:empty}

\textbf{The job.} Find the missing values (\code{nan}), work round them, or fill them (NumPy's
\code{isnan}, the \code{nan} functions such as \code{nanmean}, and \code{nan\_to\_num}). In the guide:
\href{https://numpy.org/doc/stable/user/absolute_beginners.html}{the absolute beginners guide},
``More useful array operations''.

\textbf{In simplydata.} An empty value is written \code{None} or \code{float("nan")}.
\code{count\_empty\_numbers} counts them, the summaries leave them out and say so, and
\code{fill\_empty\_numbers} fills them.

\begin{lstlisting}[style=py]
readings = sdata.number_list([4.1, None, 3.9, 4.4, None])
readings.count_empty_numbers()
readings.find_average()                  # of the 3 numbers that are there
readings.fill_empty_numbers(with_value=0)
\end{lstlisting}

\unsection{Random numbers}{from-numpy:random}

\textbf{The job.} Make random numbers from a distribution, pick at random, and get the same numbers
again with a seed (NumPy's \code{default\_rng} and its Generator methods). In the guides:
\href{https://numpy.org/doc/stable/user/absolute_beginners.html}{the absolute beginners guide},
``Generating random numbers'', and \href{https://numpy.org/doc/stable/reference/random/index.html}%
{Random sampling}.

\textbf{In simplydata.} \code{sdata.make\_random\_numbers} makes normal or uniform numbers (and whole
numbers); \code{sdata.distribution("poisson", average=3).make\_random\_numbers(how\_many=10)} makes them
from any of 28 distributions. \code{numbers.make\_random\_picks} is \code{choice}. \code{with\_seed=}
is the seed.

\begin{lstlisting}[style=py]
heights = sdata.make_random_numbers(how_many=100, average=150, standard_deviation=8, with_seed=1)
dice = sdata.make_random_numbers(how_many=10, between=[1, 6], whole_numbers=True, with_seed=1)
arrivals = sdata.distribution("poisson", average=3).make_random_numbers(how_many=10, with_seed=1)
picked = heights.make_random_picks(how_many=5, with_seed=1)
print(picked)
\end{lstlisting}

\unsection{Saving and reading numbers}{from-numpy:files}

\textbf{The job.} Write numbers to a file and read them back (NumPy's \code{save}, \code{savetxt},
\code{loadtxt} and \code{genfromtxt}). In the guide: \href{https://numpy.org/doc/stable/user/basics.io.html}%
{I/O with NumPy}.

\textbf{In simplydata.} \code{numbers.save("file.csv")} writes CSV, Excel, Parquet, JSON or a sound
file, as the end of the name says. \code{sdata.number\_list("file.csv")} reads them back, and
\code{empty\_values=[-999]} is \code{genfromtxt}'s \code{missing\_values}. A netCDF file opens as a
block with \code{sdata.number\_block("file.nc")}.

\begin{lstlisting}[style=py]
numbers = sdata.number_list([4.1, 3.9, -999, 4.4])
numbers.save("readings.csv")
again = sdata.number_list("readings.csv", empty_values=[-999])
again.count_empty_numbers()
\end{lstlisting}

\unsection{Grids as matrices}{from-numpy:matrices}

\textbf{The job.} The determinant, the inverse and the eigenvalues of a square grid, and solving the
equations grid times x = b (NumPy's \code{np.linalg}). In the guides:
\href{https://numpy.org/doc/stable/user/absolute_beginners.html}{the absolute beginners guide},
``Working with mathematical formulas'', and
\href{https://numpy.org/doc/stable/reference/routines.linalg.html}{Linear algebra}.

\textbf{In simplydata.} \code{find\_determinant}, \code{take\_inverse} and \code{find\_eigenvalues};
\code{solve\_for(b)} is \code{solve}, and on a grid that is not square it gives the least-squares answer
(\code{lstsq}) and says so. \code{exactly=True} on each gives SymPy's exact answer, as fractions. To
solve equations, \code{solve\_for} is better than an inverse. \code{multiply\_as\_matrix\_by} is the
matrix product (\code{@}, \code{np.matmul}); \code{find\_rank}, \code{find\_size} and
\code{find\_condition\_number} are \code{matrix\_rank}, \code{norm} and \code{cond}.

\begin{lstlisting}[style=py]
grid = sdata.number_grid([[2, 1], [1, 3]])
grid.find_determinant()                  # 5
grid.find_eigenvalues()
x = grid.solve_for([5, 10])              # x = [1, 3]
print(x)
grid.solve_for([1, 1], exactly=True)     # [2/5, 1/5]
grid.multiply_as_matrix_by(sdata.number_grid([[0, 1], [1, 0]]))
grid.find_rank()                         # 2
\end{lstlisting}

\unsection{Fourier transforms}{from-numpy:fourier}

\textbf{The job.} Turn readings in time into frequencies, and back (NumPy's and SciPy's \code{fft},
\code{rfft}, \code{fftfreq} and \code{ifft}). In the guide:
\href{https://docs.scipy.org/doc/scipy/tutorial/fft.html}{Discrete Fourier Transforms (scipy.fft)}.

\textbf{In simplydata.} \code{numbers.take\_fourier\_transform(time\_step=)} gives a table with each
frequency in hertz, its size and its phase: \code{rfft} and \code{rfftfreq} in one word.
\code{sdata.find\_inverse\_fourier\_transform(table)} goes back. \code{take\_cosine\_transform} is
\code{dct}. \code{find\_spectrum} gives the power at each frequency.

\begin{lstlisting}[style=py]
wave = sdata.make_wave("sine", frequency=5, time_step=0.01, how_many=200)
frequencies = wave.take_fourier_transform(time_step=0.01)
back = sdata.find_inverse_fourier_transform(frequencies)
back.count()
\end{lstlisting}

\unsection{Areas, volumes and paths over time}{from-numpy:integration}

\textbf{The job.} The area under a curve, the volume under a surface, the running area under readings,
and how a quantity changes over time from a rate (SciPy's \code{quad}, \code{dblquad},
\code{cumulative\_trapezoid} and \code{solve\_ivp}). In the guide:
\href{https://docs.scipy.org/doc/scipy/tutorial/integrate.html}{Integration (scipy.integrate)}.

\textbf{In simplydata.} \code{sdata.find\_area\_under}, \code{sdata.find\_volume\_under},
\code{numbers.turn\_into\_running\_area} and \code{sdata.find\_path\_over\_time}. The formula is written
as text, with letters. \code{sdata.find\_path\_between\_ends} is \code{solve\_bvp}, for a path with a value
given at each end, and \code{sdata.find\_area\_points} gives the points and weights of Gaussian quadrature
(\code{roots\_legendre} and its kin).

\begin{lstlisting}[style=py]
sdata.find_area_under("exp(-x**2)", between=[0, "infinity"])
sdata.find_volume_under("x * y", between={"x": [0, 1], "y": [0, 2]})
flow = sdata.number_list([0.0, 1.5, 2.0, 1.0, 0.5])
flow.turn_into_running_area(time_step=60)
path = sdata.find_path_over_time("dN/dt = r * N", at_start={"N": 10}, between=[0, 10], known={"r": 0.3})
\end{lstlisting}

\unsection{Smallest, biggest and solutions}{from-numpy:optimization}

\textbf{The job.} Find where a formula is smallest or biggest, solve equations, fit a curve to data,
and solve least squares with limits (SciPy's \code{minimize}, \code{root}, \code{curve\_fit},
\code{linprog}, \code{nnls} and \code{lsq\_linear}). In the guide:
\href{https://docs.scipy.org/doc/scipy/tutorial/optimize.html}{Optimization (scipy.optimize)}.

\textbf{In simplydata.} \code{sdata.find\_where\_smallest} and \code{sdata.find\_where\_biggest} (with
\code{rules=} for linear programming), \code{sdata.solve}, \code{table.fit\_formula} for curves, and
\code{grid.solve\_for(b, between=[0, None])} for answers that may not be negative, as in a diet mix or
the unmixing of spectra. \code{sdata.find\_best\_pairing(costs=grid)} is \code{linear\_sum\_assignment}: it
pairs each row with one column so that the total cost is smallest.

\begin{lstlisting}[style=py]
sdata.find_where_smallest("x**2 - 4*x + 1", between=[0, 5])
sdata.find_where_smallest("2*x + 3*y", rules=["x + y >= 4", "x >= 0", "y >= 0"])
sdata.solve("x**3 - 2*x = 5", between=[0, 3])
# one row per nutrient (protein, fibre, fat), one column per food: grams of the nutrient in one portion
foods = sdata.number_grid([[20, 2, 5], [10, 15, 1], [2, 8, 30]])
amounts = foods.solve_for([60, 40, 50], between=[0, None])          # portions of each food, never below 0
grid = sdata.number_grid([[4, 1, 3], [2, 0, 5], [3, 2, 2]])         # cost of each worker (row) for each job
sdata.find_best_pairing(costs=grid)
\end{lstlisting}

\unsection{Values between the points}{from-numpy:interpolation}

\textbf{The job.} Find values between measured points, on a line, on a map or in a block (SciPy's
\code{interp1d}, \code{CubicSpline}, \code{griddata} and \code{RegularGridInterpolator}). In the guide:
\href{https://docs.scipy.org/doc/scipy/tutorial/interpolate.html}{Interpolation (scipy.interpolate)}.

\textbf{In simplydata.} \code{table.find\_values\_at} on a table's columns,
\code{table.make\_map\_of} for scattered points on a map, and \code{block.find\_values\_at} inside a
labelled block. \code{sdata.curve\_through\_points(x=..., y=...)} is a spline you can keep (SciPy's
\code{make\_interp\_spline} and \code{CubicSpline}): read it anywhere with \code{curve.find\_value\_at}, and
take its slope, area and zeros with \code{find\_slope\_at}, \code{find\_area\_between} and \code{find\_zeros}.

\begin{lstlisting}[style=py]
growth = sdata.get_example("growth")
growth.keep_rows_where("country", is_equal_to="Applia")
growth.find_values_at("population", at={"year": [2001.5, 2003.5]})
curve = sdata.curve_through_points(x=[0, 5, 10, 20], y=[8.1, 7.9, 6.2, 5.0])
curve.find_value_at(12)
curve.find_slope_at(12)
\end{lstlisting}

\unsection{Signals}{from-numpy:signals}

\textbf{The job.} Smooth readings, keep or remove frequencies with a filter, find peaks, and find the
spectrum (SciPy's \code{savgol\_filter}, \code{butter} with \code{filtfilt}, \code{find\_peaks} and
\code{welch}). In the guide: \href{https://docs.scipy.org/doc/scipy/tutorial/signal.html}{Signal
Processing (scipy.signal)}.

\textbf{In simplydata.} \code{smooth}, \code{keep\_each\_frequency\_that} and
\code{remove\_each\_frequency\_that}, \code{find\_each\_peak\_that} and \code{find\_spectrum}. Each
takes \code{time\_step=}, the time between readings, so frequencies are in hertz. The filter designs of
\code{scipy.signal} (Butterworth, Chebyshev, elliptic, Bessel and FIR filters) are values of
\code{sdata.frequency\_filter}; \code{system.move\_poles\_to} is \code{place\_poles}, and
\code{find\_phase\_locking\_to} measures how steady the phase is.

\begin{lstlisting}[style=py]
signal = sdata.make_wave("sine", frequency=2, time_step=0.01, how_many=300)
noise = sdata.make_random_numbers(how_many=300, average=0, standard_deviation=0.3, with_seed=1)
signal.add(noise)
signal.keep_each_frequency_that(is_below=5, time_step=0.01)
signal.find_each_peak_that(is_at_least=0.5)
\end{lstlisting}

\unsection{More linear algebra}{from-numpy:linalg}

\textbf{The job.} Break a matrix down into parts (SVD, QR, LU, Cholesky, Schur and more), the matrix
exponential, and the null space (SciPy's \code{scipy.linalg}). In the guide:
\href{https://docs.scipy.org/doc/scipy/tutorial/linalg.html}{Linear Algebra (scipy.linalg)}.

\textbf{In simplydata.} \code{grid.break\_down(test\_method="svd")} gives the parts as a dict of
grids and prints what each part is; the other methods are \code{"qr"}, \code{"lu"},
\code{"cholesky"}, \code{"schur"}, \code{"polar"}, \code{"ldl"}, \code{"hessenberg"} and \code{"qz"}.
\code{take\_matrix\_exp} is \code{expm} and \code{find\_null\_space} is \code{null\_space}. The other
matrix functions are one family word, \code{grid.take\_matrix\_function("log")} (\code{logm}, \code{sqrtm},
\code{cosm} ...); \code{check\_if\_grid\_is("symmetric")} checks a kind of grid, and
\code{sdata.solve\_matrix\_equation} solves the Sylvester, Lyapunov and Riccati equations.

\begin{lstlisting}[style=py]
grid = sdata.number_grid([[4, 2], [2, 3]])
parts = grid.break_down(test_method="cholesky")
print(parts["U"])
parts = grid.break_down(test_method="svd")
print(parts["singular values"])
grid.find_null_space()
grid.take_matrix_function("square root")                 # SciPy's sqrtm
grid.check_if_grid_is("symmetric")
\end{lstlisting}

\unsection{Points in space}{from-numpy:spatial}

\textbf{The job.} The nearest points, distances between points, and the fence round a set of points
(SciPy's \code{KDTree}, \code{distance.cdist} and \code{ConvexHull}). In the guide:
\href{https://docs.scipy.org/doc/scipy/tutorial/spatial.html}{Spatial data structures and algorithms
(scipy.spatial)}.

\textbf{In simplydata.} These are table words, because points are rows with an east and a north:
\code{table.find\_nearest\_rows\_to}, \code{table.find\_distances\_to\_rows\_of} and
\code{table.find\_fence\_around}. \code{table.find\_triangles\_between\_points} is \code{Delaunay} and
\code{table.find\_regions\_around\_points} is \code{Voronoi}, with the area of each region.

\begin{lstlisting}[style=py]
island = sdata.get_example("island")
island.find_nearest_rows_to({"east": 4, "north": 1}, how_many=3)
island.find_fence_around(x="east", y="north")
island = sdata.get_example("island")
island.find_triangles_between_points(x="east", y="north")
\end{lstlisting}

\unsection{Special functions}{from-numpy:special}

\textbf{The job.} The gamma and beta functions, the error function, Bessel functions, and the number of
ways to choose (SciPy's \code{scipy.special}). In the guide:
\href{https://docs.scipy.org/doc/scipy/tutorial/special.html}{Special Functions (scipy.special)}.

\textbf{In simplydata.} A formula can use \code{gamma}, \code{lgamma}, \code{digamma}, \code{beta},
\code{erf}, \code{erfc} and \code{besselJ}. \code{sdata.count\_ways\_to\_choose} is \code{comb} and
\code{sdata.count\_ways\_to\_order} is \code{perm}; \code{take\_log\_odds} and \code{turn\_into\_chances}
are \code{logit} and \code{expit}. Every other function of \code{scipy.special} is a name of
\code{numbers.take\_function} (see ``Maths on every number''), and the same names work in a formula with
\code{\_} for spaces, such as \code{bessel\_j(1, x)}. \code{sdata.find\_zeros\_of("bessel J")} gives the zeros
(\code{jn\_zeros} and its kin), and \code{sdata.polynomial("legendre", degree=3)} makes the polynomials of
\code{scipy.special} and \code{np.polynomial} as formulas.

\begin{lstlisting}[style=py]
formula = sdata.formula("gamma(x)")
formula.find_value_at(x=4.5)
sdata.count_ways_to_choose(3, out_of=10)                 # 120
chances = sdata.number_list([0.1, 0.5, 0.9])
chances.take_log_odds()
sdata.find_zeros_of("bessel J", how_many=3)
formula = sdata.polynomial("legendre", degree=3)
formula.find_value_at(x=0.5)
\end{lstlisting}

\unsection{Statistics}{from-numpy:statistics}

\textbf{The job.} Use a distribution (density, chances, percentiles, random numbers), fit one to data,
and run tests (SciPy's \code{scipy.stats}). In the guide:
\href{https://docs.scipy.org/doc/scipy/tutorial/stats.html}{Statistics (scipy.stats)}.

\textbf{In simplydata.} \code{sdata.distribution(...)} with \code{find\_density\_at} (\code{pdf}),
\code{find\_chance\_that\_value} (\code{cdf} and \code{sf}) and \code{find\_percentile} (\code{ppf});
\code{numbers.fit\_distribution} is \code{fit}. The tests are table words, such as
\code{table.compare\_averages\_of}; the rest of this manual explains them. Every distribution of
\code{scipy.stats} has a plain name, such as \code{sdata.distribution("skew normal", shape=4)}.

\begin{lstlisting}[style=py]
distribution = sdata.distribution("normal", average=50, standard_deviation=10)
distribution.find_chance_that_value(is_below=60)
distribution.find_percentile(95)
waits = sdata.make_random_numbers(how_many=200, average=10, standard_deviation=2, with_seed=1)
waits.fit_distribution("normal")
\end{lstlisting}

\unsection{Images and grids of cells}{from-numpy:ndimage}

\textbf{The job.} Smooth a grid, find patches of cells that pass a test, turn and resize a grid
(SciPy's \code{scipy.ndimage}: \code{gaussian\_filter}, \code{label}, \code{rotate} and \code{zoom}).
In the guide: \href{https://docs.scipy.org/doc/scipy/tutorial/ndimage.html}{Multidimensional image
processing (scipy.ndimage)}.

\textbf{In simplydata.} \code{grid.smooth}, \code{grid.find\_each\_patch\_that}, \code{grid.rotate} and
\code{grid.change\_size\_to}; on a block, \code{block.smooth(over=[...])}. \code{grid.find\_edges()} is
\code{sobel}; \code{grid.change\_patch\_shapes("grow")} is one family word for \code{binary\_dilation},
\code{binary\_erosion}, opening, closing and filling holes; \code{find\_distance\_to\_nearest} is
\code{distance\_transform\_edt}.

\begin{lstlisting}[style=py]
grid = sdata.number_grid([[0, 0, 1, 1], [0, 0, 1, 1], [5, 0, 0, 0], [5, 5, 0, 0]])
grid.find_each_patch_that(is_above=0)
grid.change_size_to(rows=8, columns=8)
grid.rotate(by_degrees=15)
grid = sdata.number_grid([[0, 0, 0, 0], [0, 1, 1, 0], [0, 1, 1, 0], [0, 0, 0, 0]])
grid.find_edges()
\end{lstlisting}

\unsection{Sound files and other files}{from-numpy:file-io}

\textbf{The job.} Read and write sound files, and other file types (SciPy's \code{scipy.io}, such as
\code{wavfile}). In the guide: \href{https://docs.scipy.org/doc/scipy/tutorial/io.html}{File IO
(scipy.io)}.

\textbf{In simplydata.} \code{numbers.save("sound.wav", time\_step=)} and
\code{sdata.number\_list("sound.wav")}. netCDF files open as a block, and
\code{sdata.show\_whats\_inside("file.nc")} lists what is in one.

\begin{lstlisting}[style=py]
tone = sdata.make_wave("sine", frequency=440, time_step=1 / 8000, how_many=8000)
tone.times(0.5)
tone.save("tone.wav", time_step=1 / 8000)
again = sdata.number_list("tone.wav")
again.count()
\end{lstlisting}

\unsection{The jobs left to NumPy and SciPy}{from-numpy:left-out}

simplydata does the jobs most people need in research and teaching, with one word each. Many jobs that
were once left to NumPy now have words: \code{np.diff} is \code{turn\_into\_changes}, \code{np.histogram} is
\code{count\_in\_bins}, \code{np.argsort} is \code{sort(show="positions")}, \code{np.meshgrid} is
\code{sdata.make\_coordinate\_grids}, \code{np.linalg.norm} is \code{find\_size}, the matrix product is
\code{multiply\_as\_matrix\_by}, and \code{np.polynomial} is \code{sdata.polynomial}. These jobs are left out
on purpose. Each one is a hand-over: \code{array = numbers.as\_numpy()} (or \code{grid.as\_numpy()},
\code{block.as\_numpy()}), then NumPy or SciPy code, then \code{sdata.number\_list(...)},
\code{sdata.number\_grid(...)} or \code{sdata.number\_block(...)} to come back.
\begin{itemize}
  \item \textbf{Sums over named indices} (\code{np.einsum}, \code{np.linalg.tensorinv}): einsum is its own
    notation; \code{multiply\_as\_matrix\_by} covers the usual uses.
  \item \textbf{Data types, copies and views, structured arrays and strings in arrays}: simplydata keeps
    numbers as decimals and makes copies, so these questions do not come up.
  \item \textbf{Sparse arrays} (\code{scipy.sparse}) and their eigenvalue and graph routines.
  \item \textbf{The intersection of half-spaces} (\code{scipy.spatial.HalfspaceIntersection}), expert
    geometry with no data job.
  \item \textbf{Steps inside a solver} (\code{qr\_update}, \code{expm\_frechet}, the BLAS and LAPACK
    functions): \code{sdata.help} names the word that runs them.
\end{itemize}

One more example: a matrix product, and back.

\begin{lstlisting}[style=py]
a = sdata.number_grid([[1, 2], [3, 4]])
b = sdata.number_grid([[0, 1], [1, 0]])
product = a.as_numpy() @ b.as_numpy()    # NumPy's matrix product
grid = sdata.number_grid(product)        # back to simplydata
grid.find_determinant()
\end{lstlisting}
).
% The topics follow the order of the NumPy absolute beginners guide and the NumPy user guide (fundamentals), then the
% SciPy user guide, used as a checklist only: the text and the examples are our own, and each section links to the
% official page. Every example in a py box is run by tests/test_from_numpy_chapter.py (with "import simplydata as
% sdata" in front), so it cannot go out of date.

\unchapter{From NumPy and SciPy to simplydata}{from-numpy}

This chapter is for you if you already know NumPy or SciPy, or if you want to learn them after
simplydata. It follows the topics of the official guides, in their order: the
\href{https://numpy.org/doc/stable/user/absolute_beginners.html}{NumPy absolute beginners guide} and
the \href{https://numpy.org/doc/stable/user/basics.html}{NumPy fundamentals}, then the
\href{https://docs.scipy.org/doc/scipy/tutorial/index.html}{SciPy user guide}. For each topic it says
what the job is, which simplydata words do it, and a short example. Some jobs are left to NumPy and
SciPy on purpose. For those, the chapter shows how to hand the numbers over and how to bring the
answer back.

Every example starts with \code{import simplydata as sdata}. It is written once here, not in each box.

\unsection{Five things to know first}{from-numpy:first}

\begin{enumerate}
  \item \textbf{Three makers, one family.} \code{sdata.number\_list} makes a list of numbers (NumPy's
    1-D array), \code{sdata.number\_grid} makes rows and columns (a 2-D array) and
    \code{sdata.number\_block} makes numbers in any number of directions, each with a name (an n-D
    array, or an xarray DataArray). The words of numbers work on grids and blocks too.
  \item \textbf{Positions count from 1.} \code{numbers.get\_number\_at(1)} is the first number. NumPy
    counts from 0, so the sentence a word prints gives NumPy's position too.
  \item \textbf{A word changes the object, and \code{undo} takes it back.} \code{numbers.take\_log()}
    changes \code{numbers} itself and prints what it did. \code{numbers.undo()} brings back the numbers
    from before. A word that finds something (an average, a determinant) gives it back and changes
    nothing.
  \item \textbf{Empty values stay empty.} NumPy writes a missing value as \code{nan}. simplydata calls it
    an empty value, says how many there are, and leaves them out of averages and totals, as NumPy's
    \code{nanmean} does.
  \item \textbf{The door is always open.} \code{array = numbers.as\_numpy()} gives a plain NumPy array,
    for any NumPy or SciPy call. \code{sdata.number\_list(array)}, \code{sdata.number\_grid(array)} and
    \code{sdata.number\_block(array)} bring the answer back.
\end{enumerate}

\begin{lstlisting}[style=py]
numbers = sdata.number_list([2.5, 3.1, 4.8, 5.0, 6.2])
array = numbers.as_numpy()               # a NumPy array: use any NumPy or SciPy code here
import numpy as np
steps = np.diff(array)
back = sdata.number_list(steps)          # and back to simplydata
back.find_average()
\end{lstlisting}

Do you know a NumPy, SciPy or SymPy name, but not the simplydata word? Ask:
\code{sdata.help("np.mean")}, \code{sdata.help("np.linalg.solve")}, \code{sdata.help("rng.normal")}
or \code{sdata.help("scipy.stats.gamma")}. The name can have \code{np.}, \code{np.linalg.},
\code{scipy.}, \code{stats.} or \code{optimize.} in front, and the answer is the word for that package:
\code{np.mean} is the average of numbers, but a bare \code{mean} is the average of a table column.

\unsection{Making numbers, grids and blocks}{from-numpy:making}

\textbf{The job.} Make an array: from a list, as a run of evenly spaced numbers, filled with one
value, or of a given shape. In the guides:
\href{https://numpy.org/doc/stable/user/absolute_beginners.html}{the absolute beginners guide}, ``How
to create a basic array'', and \href{https://numpy.org/doc/stable/user/basics.creation.html}{Array
creation}.

\textbf{In simplydata.} \code{sdata.number\_list([...])} is \code{np.array}; \code{sdata.number\_list(0,
to=1, how\_many=11)} is \code{np.linspace} and \code{step=} is \code{np.arange}.
\code{sdata.number\_grid(rows=2, columns=3, filled\_with=0)} is \code{np.zeros} (and \code{np.full});
\code{sdata.number\_block(of\_shape=(2, 3, 4), filled\_with=1)} is \code{np.ones} with any shape.
Special grids come from \code{kind=}: \code{sdata.number\_grid(kind="identity", size=3)} is \code{np.eye},
\code{kind="diagonal"} is \code{np.diag}, and \code{"toeplitz"}, \code{"hilbert"} or \code{"pascal"} are the
grids of \code{scipy.linalg}. \code{sdata.make\_coordinate\_grids(x=..., y=...)} is \code{np.meshgrid}.

\begin{lstlisting}[style=py]
temperatures = sdata.number_list([8.3, 10.3, 11.0, 9.2, 7.5])
steps = sdata.number_list(0, to=1, how_many=11)          # 0, 0.1, ..., 1
every_five = sdata.number_list(0, to=50, step=5)
grid = sdata.number_grid([[1, 2, 3], [4, 5, 6]])
zeros = sdata.number_grid(rows=2, columns=3, filled_with=0)
block = sdata.number_block(of_shape=(2, 3, 4), filled_with=1)
print(grid)
grid = sdata.number_grid(kind="identity", size=3)       # NumPy's eye
a, b = sdata.make_coordinate_grids(x=[0, 1, 2], y=[10, 20])    # NumPy's meshgrid
\end{lstlisting}

\unsection{Shape and size}{from-numpy:shape}

\textbf{The job.} Know how many numbers there are, how many rows and columns, and how many directions
(NumPy's \code{size}, \code{shape} and \code{ndim}). In the guide:
\href{https://numpy.org/doc/stable/user/absolute_beginners.html}{the absolute beginners guide}, ``How
do you know the shape and size of an array?''.

\textbf{In simplydata.} \code{numbers.count()} is \code{size}; \code{grid.find\_shape()} and
\code{block.find\_shape()} are \code{shape}; \code{block.count\_directions()} is \code{ndim}.
\code{grid.count\_rows()} and \code{grid.count\_columns()} say one size each.

\begin{lstlisting}[style=py]
grid = sdata.number_grid([[1, 2, 3], [4, 5, 6]])
grid.find_shape()                        # (2, 3)
grid.count_rows()
block = sdata.number_block(of_shape=(2, 3, 4), filled_with=0)
block.count_directions()                 # 3
\end{lstlisting}

\unsection{Picking numbers out}{from-numpy:indexing}

\textbf{The job.} Take one number, a slice, a row or a column, or the numbers that pass a test
(NumPy's indexing and slicing, and boolean indexing). In the guides:
\href{https://numpy.org/doc/stable/user/absolute_beginners.html}{the absolute beginners guide},
``Indexing and slicing'', and \href{https://numpy.org/doc/stable/user/basics.indexing.html}{Indexing
on ndarrays}.

\textbf{In simplydata.} \code{get\_number\_at} is \code{a[i]} (from 1); \code{keep\_numbers\_by\_position(2,
to=4)} is a slice and keeps both ends; \code{grid.get\_row} and \code{grid.get\_column} are \code{a[i, :]}
and \code{a[:, j]}. \code{keep\_each\_number\_that(is\_above=5)} is \code{a[a > 5]}.
\code{grid.get\_part("diagonal")} and \code{grid.get\_part("upper triangle")} are \code{np.diagonal} and
\code{np.triu}; \code{find\_positions\_where(is\_above=5)} is \code{np.nonzero} (counting from 1).

\begin{lstlisting}[style=py]
numbers = sdata.number_list([4, 8, 15, 16, 23, 42])
numbers.get_number_at(1)                 # 4, the first number (NumPy's a[0])
numbers.keep_numbers_by_position(2, to=4)
print(numbers)                           # 8, 15, 16
numbers.undo()
numbers.keep_each_number_that(is_above=10)
grid = sdata.number_grid([[1, 2, 3], [4, 5, 6]])
grid.get_row(2)
grid.get_column(3)
grid.get_part("diagonal")                # 1 and 5
\end{lstlisting}

\unsection{Maths on every number}{from-numpy:maths}

\textbf{The job.} Add, multiply, take logs and square roots of every number at once (NumPy's universal
functions), also between two arrays of the same length. In the guides:
\href{https://numpy.org/doc/stable/user/absolute_beginners.html}{the absolute beginners guide}, ``Basic
array operations'', and \href{https://numpy.org/doc/stable/user/basics.ufuncs.html}{Universal
functions}.

\textbf{In simplydata.} \code{add}, \code{subtract}, \code{times}, \code{divide\_by},
\code{raise\_to\_power}, \code{take\_square\_root}, \code{take\_log}, \code{take\_exp},
\code{take\_sin}, \code{round\_to} and their family. Each takes one number or numbers of the same
length. \code{+}, \code{-}, \code{*} and \code{/} work too and give new numbers.
Every other function of NumPy and \code{scipy.special} is one family word:
\code{numbers.take\_function("hyperbolic sine")} is \code{np.sinh}, and
\code{take\_function("bessel J", order=1)} is \code{special.jv}. The library names work too, and
\code{numbers.take\_function()} with no name lists them all. \code{work\_on\_bits} does the bit operations
(\code{np.bitwise\_and} and the rest). \code{turn\_into\_changes} is \code{np.diff},
\code{turn\_into\_slopes} is \code{np.gradient} and \code{find\_product} is \code{np.prod}.

\begin{lstlisting}[style=py]
prices = sdata.number_list([1.20, 3.50, 2.25])
prices.times(1.25)                       # every price with 25 % added
prices.round_to(with_decimals=2)
counts = sdata.number_list([10, 100, 1000])
counts.take_log(base=10)
both = sdata.number_list([1, 2, 3]) + sdata.number_list([10, 20, 30])
print(both)
numbers = sdata.number_list([0.5, 1.0, 2.0])
numbers.take_function("hyperbolic sine")                 # NumPy's sinh
numbers.undo()
numbers.take_function("bessel J", order=1)               # SciPy's special.jv
numbers = sdata.number_list([3, 5, 9, 10])
numbers.turn_into_changes()                              # 2, 4, 1: NumPy's diff
\end{lstlisting}

\unsection{Stretching to the same shape}{from-numpy:broadcasting}

\textbf{The job.} Work on two arrays of different shapes by stretching the smaller one (NumPy's
broadcasting). In the guide: \href{https://numpy.org/doc/stable/user/basics.broadcasting.html}%
{Broadcasting}.

\textbf{In simplydata.} The maths words stretch by themselves, as NumPy does, and print a note when
they do. \code{block.show\_how\_it\_lines\_up\_with(other)} explains how two blocks line up before you
use them, and \code{sdata.stretch\_to\_same\_shape(a, b)} gives both at the same shape.
\code{sdata.hide\_stretching\_notes()} turns the notes off.

\begin{lstlisting}[style=py]
a = sdata.number_block(of_shape=(3, 4), filled_with=1)
b = sdata.number_block(of_shape=(1, 4), filled_with=2)
a.show_how_it_lines_up_with(b)
a_big, b_big = sdata.stretch_to_same_shape(a, b)
b_big.find_shape()
\end{lstlisting}

\unsection{Summaries}{from-numpy:summaries}

\textbf{The job.} The total, average, smallest and biggest number, the standard deviation and
percentiles, of all numbers or along one direction (NumPy's aggregation functions with
\code{axis=}). In the guide: \href{https://numpy.org/doc/stable/user/absolute_beginners.html}{the
absolute beginners guide}, ``More useful array operations''.

\textbf{In simplydata.} \code{find\_total}, \code{find\_average}, \code{find\_median},
\code{find\_smallest}, \code{find\_biggest}, \code{find\_standard\_deviation} and
\code{find\_percentile}. On a block, \code{find\_average(over="month")} averages along a named
direction, as \code{axis=} does with a number. The standard deviation is the sample one (NumPy's
\code{ddof=1}), as in statistics books.

\begin{lstlisting}[style=py]
numbers = sdata.number_list([8.3, 10.3, 11.0, 9.2, 7.5])
numbers.find_average()
numbers.find_standard_deviation()        # sample standard deviation (NumPy's ddof=1)
numbers.find_percentile(90)
block = sdata.number_block(of_shape=(2, 3), filled_with=1, directions_named=["year", "month"])
block.find_average(over="month")
\end{lstlisting}

\unsection{Changing the shape}{from-numpy:reshaping}

\textbf{The job.} Reshape, turn rows into columns, join arrays and split them, add or remove a
direction (NumPy's \code{reshape}, \code{transpose}, \code{concatenate}, \code{split},
\code{expand\_dims} and \code{squeeze}). In the guide:
\href{https://numpy.org/doc/stable/user/absolute_beginners.html}{the absolute beginners guide}, ``Can
you reshape an array?'' and ``Transposing and reshaping a matrix''.

\textbf{In simplydata.} \code{numbers.make\_number\_grid(rows=, columns=)} makes a grid from numbers;
\code{block.change\_shape\_to(...)} changes a block's own shape. \code{grid.swap\_rows\_and\_columns()}
is \code{.T}. \code{block.add\_numbers\_from(other, over=)} joins, \code{block.split\_into\_pieces}
splits, and \code{add\_direction} and \code{remove\_directions\_of\_size\_one} are \code{expand\_dims}
and \code{squeeze}.

\begin{lstlisting}[style=py]
numbers = sdata.number_list([1, 2, 3, 4, 5, 6])
grid = numbers.make_number_grid(rows=2, columns=3)
grid.swap_rows_and_columns()
grid.find_shape()                        # (3, 2)
block = sdata.number_block(of_shape=(2, 6), filled_with=0, directions_named=["row", "column"])
block.change_shape_to(2, 3, 2)
block.find_shape()
\end{lstlisting}

\unsection{Sorting, unique numbers and counts}{from-numpy:sorting}

\textbf{The job.} Sort, keep each different number once, and count how often each comes (NumPy's
\code{sort}, \code{unique} and \code{unique\_counts}). In the guide:
\href{https://numpy.org/doc/stable/user/absolute_beginners.html}{the absolute beginners guide},
``Adding, removing, and sorting elements'' and ``How to get unique items and counts''.

\textbf{In simplydata.} \code{sort} (\code{biggest\_first=True} for the other way),
\code{remove\_duplicate\_numbers} and \code{count\_each}, which gives a table of each number and its
count. \code{sort(show="positions")} is \code{np.argsort}, and \code{count\_in\_bins(how\_many\_bins=3)} is
\code{np.histogram} as numbers.

\begin{lstlisting}[style=py]
numbers = sdata.number_list([3, 1, 2, 3, 3, 1])
numbers.count_each()
numbers.sort(biggest_first=True)
numbers.remove_duplicate_numbers()
print(numbers)
numbers = sdata.number_list([3, 1, 2, 3, 3, 1])
numbers.count_in_bins(how_many_bins=3)
\end{lstlisting}

\unsection{Yes or no, and conditions}{from-numpy:conditions}

\textbf{The job.} Ask a question of every number and use the answers: count them, keep the numbers
where the answer is yes, or replace them (NumPy's comparisons, boolean masks, \code{where},
\code{count\_nonzero}, \code{any} and \code{all}). In the guide:
\href{https://numpy.org/doc/stable/user/basics.indexing.html}{Indexing on ndarrays}, ``Boolean array
indexing''.

\textbf{In simplydata.} \code{turn\_into\_yes\_or\_no(is\_above=)} is the mask;
\code{keep\_numbers\_where}, \code{replace\_numbers\_where} and \code{count\_numbers\_that} use it, and
\code{check\_if\_any\_number} and \code{check\_if\_all\_numbers} are \code{any} and \code{all}.

\begin{lstlisting}[style=py]
rain = sdata.number_list([0.0, 2.5, 0.0, 12.1, 4.4, 0.0])
rain.count_numbers_that(is_above=0)      # 3 days with rain
wet = rain.turn_into_yes_or_no(is_above=0)
rain.check_if_any_number(is_above=10)
rain.replace_numbers_where(wet, with_value=1)
print(rain)
\end{lstlisting}

\unsection{Empty values}{from-numpy:empty}

\textbf{The job.} Find the missing values (\code{nan}), work round them, or fill them (NumPy's
\code{isnan}, the \code{nan} functions such as \code{nanmean}, and \code{nan\_to\_num}). In the guide:
\href{https://numpy.org/doc/stable/user/absolute_beginners.html}{the absolute beginners guide},
``More useful array operations''.

\textbf{In simplydata.} An empty value is written \code{None} or \code{float("nan")}.
\code{count\_empty\_numbers} counts them, the summaries leave them out and say so, and
\code{fill\_empty\_numbers} fills them.

\begin{lstlisting}[style=py]
readings = sdata.number_list([4.1, None, 3.9, 4.4, None])
readings.count_empty_numbers()
readings.find_average()                  # of the 3 numbers that are there
readings.fill_empty_numbers(with_value=0)
\end{lstlisting}

\unsection{Random numbers}{from-numpy:random}

\textbf{The job.} Make random numbers from a distribution, pick at random, and get the same numbers
again with a seed (NumPy's \code{default\_rng} and its Generator methods). In the guides:
\href{https://numpy.org/doc/stable/user/absolute_beginners.html}{the absolute beginners guide},
``Generating random numbers'', and \href{https://numpy.org/doc/stable/reference/random/index.html}%
{Random sampling}.

\textbf{In simplydata.} \code{sdata.make\_random\_numbers} makes normal or uniform numbers (and whole
numbers); \code{sdata.distribution("poisson", average=3).make\_random\_numbers(how\_many=10)} makes them
from any of 28 distributions. \code{numbers.make\_random\_picks} is \code{choice}. \code{with\_seed=}
is the seed.

\begin{lstlisting}[style=py]
heights = sdata.make_random_numbers(how_many=100, average=150, standard_deviation=8, with_seed=1)
dice = sdata.make_random_numbers(how_many=10, between=[1, 6], whole_numbers=True, with_seed=1)
arrivals = sdata.distribution("poisson", average=3).make_random_numbers(how_many=10, with_seed=1)
picked = heights.make_random_picks(how_many=5, with_seed=1)
print(picked)
\end{lstlisting}

\unsection{Saving and reading numbers}{from-numpy:files}

\textbf{The job.} Write numbers to a file and read them back (NumPy's \code{save}, \code{savetxt},
\code{loadtxt} and \code{genfromtxt}). In the guide: \href{https://numpy.org/doc/stable/user/basics.io.html}%
{I/O with NumPy}.

\textbf{In simplydata.} \code{numbers.save("file.csv")} writes CSV, Excel, Parquet, JSON or a sound
file, as the end of the name says. \code{sdata.number\_list("file.csv")} reads them back, and
\code{empty\_values=[-999]} is \code{genfromtxt}'s \code{missing\_values}. A netCDF file opens as a
block with \code{sdata.number\_block("file.nc")}.

\begin{lstlisting}[style=py]
numbers = sdata.number_list([4.1, 3.9, -999, 4.4])
numbers.save("readings.csv")
again = sdata.number_list("readings.csv", empty_values=[-999])
again.count_empty_numbers()
\end{lstlisting}

\unsection{Grids as matrices}{from-numpy:matrices}

\textbf{The job.} The determinant, the inverse and the eigenvalues of a square grid, and solving the
equations grid times x = b (NumPy's \code{np.linalg}). In the guides:
\href{https://numpy.org/doc/stable/user/absolute_beginners.html}{the absolute beginners guide},
``Working with mathematical formulas'', and
\href{https://numpy.org/doc/stable/reference/routines.linalg.html}{Linear algebra}.

\textbf{In simplydata.} \code{find\_determinant}, \code{take\_inverse} and \code{find\_eigenvalues};
\code{solve\_for(b)} is \code{solve}, and on a grid that is not square it gives the least-squares answer
(\code{lstsq}) and says so. \code{exactly=True} on each gives SymPy's exact answer, as fractions. To
solve equations, \code{solve\_for} is better than an inverse. \code{multiply\_as\_matrix\_by} is the
matrix product (\code{@}, \code{np.matmul}); \code{find\_rank}, \code{find\_size} and
\code{find\_condition\_number} are \code{matrix\_rank}, \code{norm} and \code{cond}.

\begin{lstlisting}[style=py]
grid = sdata.number_grid([[2, 1], [1, 3]])
grid.find_determinant()                  # 5
grid.find_eigenvalues()
x = grid.solve_for([5, 10])              # x = [1, 3]
print(x)
grid.solve_for([1, 1], exactly=True)     # [2/5, 1/5]
grid.multiply_as_matrix_by(sdata.number_grid([[0, 1], [1, 0]]))
grid.find_rank()                         # 2
\end{lstlisting}

\unsection{Fourier transforms}{from-numpy:fourier}

\textbf{The job.} Turn readings in time into frequencies, and back (NumPy's and SciPy's \code{fft},
\code{rfft}, \code{fftfreq} and \code{ifft}). In the guide:
\href{https://docs.scipy.org/doc/scipy/tutorial/fft.html}{Discrete Fourier Transforms (scipy.fft)}.

\textbf{In simplydata.} \code{numbers.take\_fourier\_transform(time\_step=)} gives a table with each
frequency in hertz, its size and its phase: \code{rfft} and \code{rfftfreq} in one word.
\code{sdata.find\_inverse\_fourier\_transform(table)} goes back. \code{take\_cosine\_transform} is
\code{dct}. \code{find\_spectrum} gives the power at each frequency.

\begin{lstlisting}[style=py]
wave = sdata.make_wave("sine", frequency=5, time_step=0.01, how_many=200)
frequencies = wave.take_fourier_transform(time_step=0.01)
back = sdata.find_inverse_fourier_transform(frequencies)
back.count()
\end{lstlisting}

\unsection{Areas, volumes and paths over time}{from-numpy:integration}

\textbf{The job.} The area under a curve, the volume under a surface, the running area under readings,
and how a quantity changes over time from a rate (SciPy's \code{quad}, \code{dblquad},
\code{cumulative\_trapezoid} and \code{solve\_ivp}). In the guide:
\href{https://docs.scipy.org/doc/scipy/tutorial/integrate.html}{Integration (scipy.integrate)}.

\textbf{In simplydata.} \code{sdata.find\_area\_under}, \code{sdata.find\_volume\_under},
\code{numbers.turn\_into\_running\_area} and \code{sdata.find\_path\_over\_time}. The formula is written
as text, with letters. \code{sdata.find\_path\_between\_ends} is \code{solve\_bvp}, for a path with a value
given at each end, and \code{sdata.find\_area\_points} gives the points and weights of Gaussian quadrature
(\code{roots\_legendre} and its kin).

\begin{lstlisting}[style=py]
sdata.find_area_under("exp(-x**2)", between=[0, "infinity"])
sdata.find_volume_under("x * y", between={"x": [0, 1], "y": [0, 2]})
flow = sdata.number_list([0.0, 1.5, 2.0, 1.0, 0.5])
flow.turn_into_running_area(time_step=60)
path = sdata.find_path_over_time("dN/dt = r * N", at_start={"N": 10}, between=[0, 10], known={"r": 0.3})
\end{lstlisting}

\unsection{Smallest, biggest and solutions}{from-numpy:optimization}

\textbf{The job.} Find where a formula is smallest or biggest, solve equations, fit a curve to data,
and solve least squares with limits (SciPy's \code{minimize}, \code{root}, \code{curve\_fit},
\code{linprog}, \code{nnls} and \code{lsq\_linear}). In the guide:
\href{https://docs.scipy.org/doc/scipy/tutorial/optimize.html}{Optimization (scipy.optimize)}.

\textbf{In simplydata.} \code{sdata.find\_where\_smallest} and \code{sdata.find\_where\_biggest} (with
\code{rules=} for linear programming), \code{sdata.solve}, \code{table.fit\_formula} for curves, and
\code{grid.solve\_for(b, between=[0, None])} for answers that may not be negative, as in a diet mix or
the unmixing of spectra. \code{sdata.find\_best\_pairing(costs=grid)} is \code{linear\_sum\_assignment}: it
pairs each row with one column so that the total cost is smallest.

\begin{lstlisting}[style=py]
sdata.find_where_smallest("x**2 - 4*x + 1", between=[0, 5])
sdata.find_where_smallest("2*x + 3*y", rules=["x + y >= 4", "x >= 0", "y >= 0"])
sdata.solve("x**3 - 2*x = 5", between=[0, 3])
# one row per nutrient (protein, fibre, fat), one column per food: grams of the nutrient in one portion
foods = sdata.number_grid([[20, 2, 5], [10, 15, 1], [2, 8, 30]])
amounts = foods.solve_for([60, 40, 50], between=[0, None])          # portions of each food, never below 0
grid = sdata.number_grid([[4, 1, 3], [2, 0, 5], [3, 2, 2]])         # cost of each worker (row) for each job
sdata.find_best_pairing(costs=grid)
\end{lstlisting}

\unsection{Values between the points}{from-numpy:interpolation}

\textbf{The job.} Find values between measured points, on a line, on a map or in a block (SciPy's
\code{interp1d}, \code{CubicSpline}, \code{griddata} and \code{RegularGridInterpolator}). In the guide:
\href{https://docs.scipy.org/doc/scipy/tutorial/interpolate.html}{Interpolation (scipy.interpolate)}.

\textbf{In simplydata.} \code{table.find\_values\_at} on a table's columns,
\code{table.make\_map\_of} for scattered points on a map, and \code{block.find\_values\_at} inside a
labelled block. \code{sdata.curve\_through\_points(x=..., y=...)} is a spline you can keep (SciPy's
\code{make\_interp\_spline} and \code{CubicSpline}): read it anywhere with \code{curve.find\_value\_at}, and
take its slope, area and zeros with \code{find\_slope\_at}, \code{find\_area\_between} and \code{find\_zeros}.

\begin{lstlisting}[style=py]
growth = sdata.get_example("growth")
growth.keep_rows_where("country", is_equal_to="Applia")
growth.find_values_at("population", at={"year": [2001.5, 2003.5]})
curve = sdata.curve_through_points(x=[0, 5, 10, 20], y=[8.1, 7.9, 6.2, 5.0])
curve.find_value_at(12)
curve.find_slope_at(12)
\end{lstlisting}

\unsection{Signals}{from-numpy:signals}

\textbf{The job.} Smooth readings, keep or remove frequencies with a filter, find peaks, and find the
spectrum (SciPy's \code{savgol\_filter}, \code{butter} with \code{filtfilt}, \code{find\_peaks} and
\code{welch}). In the guide: \href{https://docs.scipy.org/doc/scipy/tutorial/signal.html}{Signal
Processing (scipy.signal)}.

\textbf{In simplydata.} \code{smooth}, \code{keep\_each\_frequency\_that} and
\code{remove\_each\_frequency\_that}, \code{find\_each\_peak\_that} and \code{find\_spectrum}. Each
takes \code{time\_step=}, the time between readings, so frequencies are in hertz. The filter designs of
\code{scipy.signal} (Butterworth, Chebyshev, elliptic, Bessel and FIR filters) are values of
\code{sdata.frequency\_filter}; \code{system.move\_poles\_to} is \code{place\_poles}, and
\code{find\_phase\_locking\_to} measures how steady the phase is.

\begin{lstlisting}[style=py]
signal = sdata.make_wave("sine", frequency=2, time_step=0.01, how_many=300)
noise = sdata.make_random_numbers(how_many=300, average=0, standard_deviation=0.3, with_seed=1)
signal.add(noise)
signal.keep_each_frequency_that(is_below=5, time_step=0.01)
signal.find_each_peak_that(is_at_least=0.5)
\end{lstlisting}

\unsection{More linear algebra}{from-numpy:linalg}

\textbf{The job.} Break a matrix down into parts (SVD, QR, LU, Cholesky, Schur and more), the matrix
exponential, and the null space (SciPy's \code{scipy.linalg}). In the guide:
\href{https://docs.scipy.org/doc/scipy/tutorial/linalg.html}{Linear Algebra (scipy.linalg)}.

\textbf{In simplydata.} \code{grid.break\_down(test\_method="svd")} gives the parts as a dict of
grids and prints what each part is; the other methods are \code{"qr"}, \code{"lu"},
\code{"cholesky"}, \code{"schur"}, \code{"polar"}, \code{"ldl"}, \code{"hessenberg"} and \code{"qz"}.
\code{take\_matrix\_exp} is \code{expm} and \code{find\_null\_space} is \code{null\_space}. The other
matrix functions are one family word, \code{grid.take\_matrix\_function("log")} (\code{logm}, \code{sqrtm},
\code{cosm} ...); \code{check\_if\_grid\_is("symmetric")} checks a kind of grid, and
\code{sdata.solve\_matrix\_equation} solves the Sylvester, Lyapunov and Riccati equations.

\begin{lstlisting}[style=py]
grid = sdata.number_grid([[4, 2], [2, 3]])
parts = grid.break_down(test_method="cholesky")
print(parts["U"])
parts = grid.break_down(test_method="svd")
print(parts["singular values"])
grid.find_null_space()
grid.take_matrix_function("square root")                 # SciPy's sqrtm
grid.check_if_grid_is("symmetric")
\end{lstlisting}

\unsection{Points in space}{from-numpy:spatial}

\textbf{The job.} The nearest points, distances between points, and the fence round a set of points
(SciPy's \code{KDTree}, \code{distance.cdist} and \code{ConvexHull}). In the guide:
\href{https://docs.scipy.org/doc/scipy/tutorial/spatial.html}{Spatial data structures and algorithms
(scipy.spatial)}.

\textbf{In simplydata.} These are table words, because points are rows with an east and a north:
\code{table.find\_nearest\_rows\_to}, \code{table.find\_distances\_to\_rows\_of} and
\code{table.find\_fence\_around}. \code{table.find\_triangles\_between\_points} is \code{Delaunay} and
\code{table.find\_regions\_around\_points} is \code{Voronoi}, with the area of each region.

\begin{lstlisting}[style=py]
island = sdata.get_example("island")
island.find_nearest_rows_to({"east": 4, "north": 1}, how_many=3)
island.find_fence_around(x="east", y="north")
island = sdata.get_example("island")
island.find_triangles_between_points(x="east", y="north")
\end{lstlisting}

\unsection{Special functions}{from-numpy:special}

\textbf{The job.} The gamma and beta functions, the error function, Bessel functions, and the number of
ways to choose (SciPy's \code{scipy.special}). In the guide:
\href{https://docs.scipy.org/doc/scipy/tutorial/special.html}{Special Functions (scipy.special)}.

\textbf{In simplydata.} A formula can use \code{gamma}, \code{lgamma}, \code{digamma}, \code{beta},
\code{erf}, \code{erfc} and \code{besselJ}. \code{sdata.count\_ways\_to\_choose} is \code{comb} and
\code{sdata.count\_ways\_to\_order} is \code{perm}; \code{take\_log\_odds} and \code{turn\_into\_chances}
are \code{logit} and \code{expit}. Every other function of \code{scipy.special} is a name of
\code{numbers.take\_function} (see ``Maths on every number''), and the same names work in a formula with
\code{\_} for spaces, such as \code{bessel\_j(1, x)}. \code{sdata.find\_zeros\_of("bessel J")} gives the zeros
(\code{jn\_zeros} and its kin), and \code{sdata.polynomial("legendre", degree=3)} makes the polynomials of
\code{scipy.special} and \code{np.polynomial} as formulas.

\begin{lstlisting}[style=py]
formula = sdata.formula("gamma(x)")
formula.find_value_at(x=4.5)
sdata.count_ways_to_choose(3, out_of=10)                 # 120
chances = sdata.number_list([0.1, 0.5, 0.9])
chances.take_log_odds()
sdata.find_zeros_of("bessel J", how_many=3)
formula = sdata.polynomial("legendre", degree=3)
formula.find_value_at(x=0.5)
\end{lstlisting}

\unsection{Statistics}{from-numpy:statistics}

\textbf{The job.} Use a distribution (density, chances, percentiles, random numbers), fit one to data,
and run tests (SciPy's \code{scipy.stats}). In the guide:
\href{https://docs.scipy.org/doc/scipy/tutorial/stats.html}{Statistics (scipy.stats)}.

\textbf{In simplydata.} \code{sdata.distribution(...)} with \code{find\_density\_at} (\code{pdf}),
\code{find\_chance\_that\_value} (\code{cdf} and \code{sf}) and \code{find\_percentile} (\code{ppf});
\code{numbers.fit\_distribution} is \code{fit}. The tests are table words, such as
\code{table.compare\_averages\_of}; the rest of this manual explains them. Every distribution of
\code{scipy.stats} has a plain name, such as \code{sdata.distribution("skew normal", shape=4)}.

\begin{lstlisting}[style=py]
distribution = sdata.distribution("normal", average=50, standard_deviation=10)
distribution.find_chance_that_value(is_below=60)
distribution.find_percentile(95)
waits = sdata.make_random_numbers(how_many=200, average=10, standard_deviation=2, with_seed=1)
waits.fit_distribution("normal")
\end{lstlisting}

\unsection{Images and grids of cells}{from-numpy:ndimage}

\textbf{The job.} Smooth a grid, find patches of cells that pass a test, turn and resize a grid
(SciPy's \code{scipy.ndimage}: \code{gaussian\_filter}, \code{label}, \code{rotate} and \code{zoom}).
In the guide: \href{https://docs.scipy.org/doc/scipy/tutorial/ndimage.html}{Multidimensional image
processing (scipy.ndimage)}.

\textbf{In simplydata.} \code{grid.smooth}, \code{grid.find\_each\_patch\_that}, \code{grid.rotate} and
\code{grid.change\_size\_to}; on a block, \code{block.smooth(over=[...])}. \code{grid.find\_edges()} is
\code{sobel}; \code{grid.change\_patch\_shapes("grow")} is one family word for \code{binary\_dilation},
\code{binary\_erosion}, opening, closing and filling holes; \code{find\_distance\_to\_nearest} is
\code{distance\_transform\_edt}.

\begin{lstlisting}[style=py]
grid = sdata.number_grid([[0, 0, 1, 1], [0, 0, 1, 1], [5, 0, 0, 0], [5, 5, 0, 0]])
grid.find_each_patch_that(is_above=0)
grid.change_size_to(rows=8, columns=8)
grid.rotate(by_degrees=15)
grid = sdata.number_grid([[0, 0, 0, 0], [0, 1, 1, 0], [0, 1, 1, 0], [0, 0, 0, 0]])
grid.find_edges()
\end{lstlisting}

\unsection{Sound files and other files}{from-numpy:file-io}

\textbf{The job.} Read and write sound files, and other file types (SciPy's \code{scipy.io}, such as
\code{wavfile}). In the guide: \href{https://docs.scipy.org/doc/scipy/tutorial/io.html}{File IO
(scipy.io)}.

\textbf{In simplydata.} \code{numbers.save("sound.wav", time\_step=)} and
\code{sdata.number\_list("sound.wav")}. netCDF files open as a block, and
\code{sdata.show\_whats\_inside("file.nc")} lists what is in one.

\begin{lstlisting}[style=py]
tone = sdata.make_wave("sine", frequency=440, time_step=1 / 8000, how_many=8000)
tone.times(0.5)
tone.save("tone.wav", time_step=1 / 8000)
again = sdata.number_list("tone.wav")
again.count()
\end{lstlisting}

\unsection{The jobs left to NumPy and SciPy}{from-numpy:left-out}

simplydata does the jobs most people need in research and teaching, with one word each. Many jobs that
were once left to NumPy now have words: \code{np.diff} is \code{turn\_into\_changes}, \code{np.histogram} is
\code{count\_in\_bins}, \code{np.argsort} is \code{sort(show="positions")}, \code{np.meshgrid} is
\code{sdata.make\_coordinate\_grids}, \code{np.linalg.norm} is \code{find\_size}, the matrix product is
\code{multiply\_as\_matrix\_by}, and \code{np.polynomial} is \code{sdata.polynomial}. These jobs are left out
on purpose. Each one is a hand-over: \code{array = numbers.as\_numpy()} (or \code{grid.as\_numpy()},
\code{block.as\_numpy()}), then NumPy or SciPy code, then \code{sdata.number\_list(...)},
\code{sdata.number\_grid(...)} or \code{sdata.number\_block(...)} to come back.
\begin{itemize}
  \item \textbf{Sums over named indices} (\code{np.einsum}, \code{np.linalg.tensorinv}): einsum is its own
    notation; \code{multiply\_as\_matrix\_by} covers the usual uses.
  \item \textbf{Data types, copies and views, structured arrays and strings in arrays}: simplydata keeps
    numbers as decimals and makes copies, so these questions do not come up.
  \item \textbf{Sparse arrays} (\code{scipy.sparse}) and their eigenvalue and graph routines.
  \item \textbf{The intersection of half-spaces} (\code{scipy.spatial.HalfspaceIntersection}), expert
    geometry with no data job.
  \item \textbf{Steps inside a solver} (\code{qr\_update}, \code{expm\_frechet}, the BLAS and LAPACK
    functions): \code{sdata.help} names the word that runs them.
\end{itemize}

One more example: a matrix product, and back.

\begin{lstlisting}[style=py]
a = sdata.number_grid([[1, 2], [3, 4]])
b = sdata.number_grid([[0, 1], [1, 0]])
product = a.as_numpy() @ b.as_numpy()    # NumPy's matrix product
grid = sdata.number_grid(product)        # back to simplydata
grid.find_determinant()
\end{lstlisting}
).
% The topics follow the order of the NumPy absolute beginners guide and the NumPy user guide (fundamentals), then the
% SciPy user guide, used as a checklist only: the text and the examples are our own, and each section links to the
% official page. Every example in a py box is run by tests/test_from_numpy_chapter.py (with "import simplydata as
% sdata" in front), so it cannot go out of date.

\unchapter{From NumPy and SciPy to simplydata}{from-numpy}

This chapter is for you if you already know NumPy or SciPy, or if you want to learn them after
simplydata. It follows the topics of the official guides, in their order: the
\href{https://numpy.org/doc/stable/user/absolute_beginners.html}{NumPy absolute beginners guide} and
the \href{https://numpy.org/doc/stable/user/basics.html}{NumPy fundamentals}, then the
\href{https://docs.scipy.org/doc/scipy/tutorial/index.html}{SciPy user guide}. For each topic it says
what the job is, which simplydata words do it, and a short example. Some jobs are left to NumPy and
SciPy on purpose. For those, the chapter shows how to hand the numbers over and how to bring the
answer back.

Every example starts with \code{import simplydata as sdata}. It is written once here, not in each box.

\unsection{Five things to know first}{from-numpy:first}

\begin{enumerate}
  \item \textbf{Three makers, one family.} \code{sdata.number\_list} makes a list of numbers (NumPy's
    1-D array), \code{sdata.number\_grid} makes rows and columns (a 2-D array) and
    \code{sdata.number\_block} makes numbers in any number of directions, each with a name (an n-D
    array, or an xarray DataArray). The words of numbers work on grids and blocks too.
  \item \textbf{Positions count from 1.} \code{numbers.get\_number\_at(1)} is the first number. NumPy
    counts from 0, so the sentence a word prints gives NumPy's position too.
  \item \textbf{A word changes the object, and \code{undo} takes it back.} \code{numbers.take\_log()}
    changes \code{numbers} itself and prints what it did. \code{numbers.undo()} brings back the numbers
    from before. A word that finds something (an average, a determinant) gives it back and changes
    nothing.
  \item \textbf{Empty values stay empty.} NumPy writes a missing value as \code{nan}. simplydata calls it
    an empty value, says how many there are, and leaves them out of averages and totals, as NumPy's
    \code{nanmean} does.
  \item \textbf{The door is always open.} \code{array = numbers.as\_numpy()} gives a plain NumPy array,
    for any NumPy or SciPy call. \code{sdata.number\_list(array)}, \code{sdata.number\_grid(array)} and
    \code{sdata.number\_block(array)} bring the answer back.
\end{enumerate}

\begin{lstlisting}[style=py]
numbers = sdata.number_list([2.5, 3.1, 4.8, 5.0, 6.2])
array = numbers.as_numpy()               # a NumPy array: use any NumPy or SciPy code here
import numpy as np
steps = np.diff(array)
back = sdata.number_list(steps)          # and back to simplydata
back.find_average()
\end{lstlisting}

Do you know a NumPy, SciPy or SymPy name, but not the simplydata word? Ask:
\code{sdata.help("np.mean")}, \code{sdata.help("np.linalg.solve")}, \code{sdata.help("rng.normal")}
or \code{sdata.help("scipy.stats.gamma")}. The name can have \code{np.}, \code{np.linalg.},
\code{scipy.}, \code{stats.} or \code{optimize.} in front, and the answer is the word for that package:
\code{np.mean} is the average of numbers, but a bare \code{mean} is the average of a table column.

\unsection{Making numbers, grids and blocks}{from-numpy:making}

\textbf{The job.} Make an array: from a list, as a run of evenly spaced numbers, filled with one
value, or of a given shape. In the guides:
\href{https://numpy.org/doc/stable/user/absolute_beginners.html}{the absolute beginners guide}, ``How
to create a basic array'', and \href{https://numpy.org/doc/stable/user/basics.creation.html}{Array
creation}.

\textbf{In simplydata.} \code{sdata.number\_list([...])} is \code{np.array}; \code{sdata.number\_list(0,
to=1, how\_many=11)} is \code{np.linspace} and \code{step=} is \code{np.arange}.
\code{sdata.number\_grid(rows=2, columns=3, filled\_with=0)} is \code{np.zeros} (and \code{np.full});
\code{sdata.number\_block(of\_shape=(2, 3, 4), filled\_with=1)} is \code{np.ones} with any shape.
Special grids come from \code{kind=}: \code{sdata.number\_grid(kind="identity", size=3)} is \code{np.eye},
\code{kind="diagonal"} is \code{np.diag}, and \code{"toeplitz"}, \code{"hilbert"} or \code{"pascal"} are the
grids of \code{scipy.linalg}. \code{sdata.make\_coordinate\_grids(x=..., y=...)} is \code{np.meshgrid}.

\begin{lstlisting}[style=py]
temperatures = sdata.number_list([8.3, 10.3, 11.0, 9.2, 7.5])
steps = sdata.number_list(0, to=1, how_many=11)          # 0, 0.1, ..., 1
every_five = sdata.number_list(0, to=50, step=5)
grid = sdata.number_grid([[1, 2, 3], [4, 5, 6]])
zeros = sdata.number_grid(rows=2, columns=3, filled_with=0)
block = sdata.number_block(of_shape=(2, 3, 4), filled_with=1)
print(grid)
grid = sdata.number_grid(kind="identity", size=3)       # NumPy's eye
a, b = sdata.make_coordinate_grids(x=[0, 1, 2], y=[10, 20])    # NumPy's meshgrid
\end{lstlisting}

\unsection{Shape and size}{from-numpy:shape}

\textbf{The job.} Know how many numbers there are, how many rows and columns, and how many directions
(NumPy's \code{size}, \code{shape} and \code{ndim}). In the guide:
\href{https://numpy.org/doc/stable/user/absolute_beginners.html}{the absolute beginners guide}, ``How
do you know the shape and size of an array?''.

\textbf{In simplydata.} \code{numbers.count()} is \code{size}; \code{grid.find\_shape()} and
\code{block.find\_shape()} are \code{shape}; \code{block.count\_directions()} is \code{ndim}.
\code{grid.count\_rows()} and \code{grid.count\_columns()} say one size each.

\begin{lstlisting}[style=py]
grid = sdata.number_grid([[1, 2, 3], [4, 5, 6]])
grid.find_shape()                        # (2, 3)
grid.count_rows()
block = sdata.number_block(of_shape=(2, 3, 4), filled_with=0)
block.count_directions()                 # 3
\end{lstlisting}

\unsection{Picking numbers out}{from-numpy:indexing}

\textbf{The job.} Take one number, a slice, a row or a column, or the numbers that pass a test
(NumPy's indexing and slicing, and boolean indexing). In the guides:
\href{https://numpy.org/doc/stable/user/absolute_beginners.html}{the absolute beginners guide},
``Indexing and slicing'', and \href{https://numpy.org/doc/stable/user/basics.indexing.html}{Indexing
on ndarrays}.

\textbf{In simplydata.} \code{get\_number\_at} is \code{a[i]} (from 1); \code{keep\_numbers\_by\_position(2,
to=4)} is a slice and keeps both ends; \code{grid.get\_row} and \code{grid.get\_column} are \code{a[i, :]}
and \code{a[:, j]}. \code{keep\_each\_number\_that(is\_above=5)} is \code{a[a > 5]}.
\code{grid.get\_part("diagonal")} and \code{grid.get\_part("upper triangle")} are \code{np.diagonal} and
\code{np.triu}; \code{find\_positions\_where(is\_above=5)} is \code{np.nonzero} (counting from 1).

\begin{lstlisting}[style=py]
numbers = sdata.number_list([4, 8, 15, 16, 23, 42])
numbers.get_number_at(1)                 # 4, the first number (NumPy's a[0])
numbers.keep_numbers_by_position(2, to=4)
print(numbers)                           # 8, 15, 16
numbers.undo()
numbers.keep_each_number_that(is_above=10)
grid = sdata.number_grid([[1, 2, 3], [4, 5, 6]])
grid.get_row(2)
grid.get_column(3)
grid.get_part("diagonal")                # 1 and 5
\end{lstlisting}

\unsection{Maths on every number}{from-numpy:maths}

\textbf{The job.} Add, multiply, take logs and square roots of every number at once (NumPy's universal
functions), also between two arrays of the same length. In the guides:
\href{https://numpy.org/doc/stable/user/absolute_beginners.html}{the absolute beginners guide}, ``Basic
array operations'', and \href{https://numpy.org/doc/stable/user/basics.ufuncs.html}{Universal
functions}.

\textbf{In simplydata.} \code{add}, \code{subtract}, \code{times}, \code{divide\_by},
\code{raise\_to\_power}, \code{take\_square\_root}, \code{take\_log}, \code{take\_exp},
\code{take\_sin}, \code{round\_to} and their family. Each takes one number or numbers of the same
length. \code{+}, \code{-}, \code{*} and \code{/} work too and give new numbers.
Every other function of NumPy and \code{scipy.special} is one family word:
\code{numbers.take\_function("hyperbolic sine")} is \code{np.sinh}, and
\code{take\_function("bessel J", order=1)} is \code{special.jv}. The library names work too, and
\code{numbers.take\_function()} with no name lists them all. \code{work\_on\_bits} does the bit operations
(\code{np.bitwise\_and} and the rest). \code{turn\_into\_changes} is \code{np.diff},
\code{turn\_into\_slopes} is \code{np.gradient} and \code{find\_product} is \code{np.prod}.

\begin{lstlisting}[style=py]
prices = sdata.number_list([1.20, 3.50, 2.25])
prices.times(1.25)                       # every price with 25 % added
prices.round_to(with_decimals=2)
counts = sdata.number_list([10, 100, 1000])
counts.take_log(base=10)
both = sdata.number_list([1, 2, 3]) + sdata.number_list([10, 20, 30])
print(both)
numbers = sdata.number_list([0.5, 1.0, 2.0])
numbers.take_function("hyperbolic sine")                 # NumPy's sinh
numbers.undo()
numbers.take_function("bessel J", order=1)               # SciPy's special.jv
numbers = sdata.number_list([3, 5, 9, 10])
numbers.turn_into_changes()                              # 2, 4, 1: NumPy's diff
\end{lstlisting}

\unsection{Stretching to the same shape}{from-numpy:broadcasting}

\textbf{The job.} Work on two arrays of different shapes by stretching the smaller one (NumPy's
broadcasting). In the guide: \href{https://numpy.org/doc/stable/user/basics.broadcasting.html}%
{Broadcasting}.

\textbf{In simplydata.} The maths words stretch by themselves, as NumPy does, and print a note when
they do. \code{block.show\_how\_it\_lines\_up\_with(other)} explains how two blocks line up before you
use them, and \code{sdata.stretch\_to\_same\_shape(a, b)} gives both at the same shape.
\code{sdata.hide\_stretching\_notes()} turns the notes off.

\begin{lstlisting}[style=py]
a = sdata.number_block(of_shape=(3, 4), filled_with=1)
b = sdata.number_block(of_shape=(1, 4), filled_with=2)
a.show_how_it_lines_up_with(b)
a_big, b_big = sdata.stretch_to_same_shape(a, b)
b_big.find_shape()
\end{lstlisting}

\unsection{Summaries}{from-numpy:summaries}

\textbf{The job.} The total, average, smallest and biggest number, the standard deviation and
percentiles, of all numbers or along one direction (NumPy's aggregation functions with
\code{axis=}). In the guide: \href{https://numpy.org/doc/stable/user/absolute_beginners.html}{the
absolute beginners guide}, ``More useful array operations''.

\textbf{In simplydata.} \code{find\_total}, \code{find\_average}, \code{find\_median},
\code{find\_smallest}, \code{find\_biggest}, \code{find\_standard\_deviation} and
\code{find\_percentile}. On a block, \code{find\_average(over="month")} averages along a named
direction, as \code{axis=} does with a number. The standard deviation is the sample one (NumPy's
\code{ddof=1}), as in statistics books.

\begin{lstlisting}[style=py]
numbers = sdata.number_list([8.3, 10.3, 11.0, 9.2, 7.5])
numbers.find_average()
numbers.find_standard_deviation()        # sample standard deviation (NumPy's ddof=1)
numbers.find_percentile(90)
block = sdata.number_block(of_shape=(2, 3), filled_with=1, directions_named=["year", "month"])
block.find_average(over="month")
\end{lstlisting}

\unsection{Changing the shape}{from-numpy:reshaping}

\textbf{The job.} Reshape, turn rows into columns, join arrays and split them, add or remove a
direction (NumPy's \code{reshape}, \code{transpose}, \code{concatenate}, \code{split},
\code{expand\_dims} and \code{squeeze}). In the guide:
\href{https://numpy.org/doc/stable/user/absolute_beginners.html}{the absolute beginners guide}, ``Can
you reshape an array?'' and ``Transposing and reshaping a matrix''.

\textbf{In simplydata.} \code{numbers.make\_number\_grid(rows=, columns=)} makes a grid from numbers;
\code{block.change\_shape\_to(...)} changes a block's own shape. \code{grid.swap\_rows\_and\_columns()}
is \code{.T}. \code{block.add\_numbers\_from(other, over=)} joins, \code{block.split\_into\_pieces}
splits, and \code{add\_direction} and \code{remove\_directions\_of\_size\_one} are \code{expand\_dims}
and \code{squeeze}.

\begin{lstlisting}[style=py]
numbers = sdata.number_list([1, 2, 3, 4, 5, 6])
grid = numbers.make_number_grid(rows=2, columns=3)
grid.swap_rows_and_columns()
grid.find_shape()                        # (3, 2)
block = sdata.number_block(of_shape=(2, 6), filled_with=0, directions_named=["row", "column"])
block.change_shape_to(2, 3, 2)
block.find_shape()
\end{lstlisting}

\unsection{Sorting, unique numbers and counts}{from-numpy:sorting}

\textbf{The job.} Sort, keep each different number once, and count how often each comes (NumPy's
\code{sort}, \code{unique} and \code{unique\_counts}). In the guide:
\href{https://numpy.org/doc/stable/user/absolute_beginners.html}{the absolute beginners guide},
``Adding, removing, and sorting elements'' and ``How to get unique items and counts''.

\textbf{In simplydata.} \code{sort} (\code{biggest\_first=True} for the other way),
\code{remove\_duplicate\_numbers} and \code{count\_each}, which gives a table of each number and its
count. \code{sort(show="positions")} is \code{np.argsort}, and \code{count\_in\_bins(how\_many\_bins=3)} is
\code{np.histogram} as numbers.

\begin{lstlisting}[style=py]
numbers = sdata.number_list([3, 1, 2, 3, 3, 1])
numbers.count_each()
numbers.sort(biggest_first=True)
numbers.remove_duplicate_numbers()
print(numbers)
numbers = sdata.number_list([3, 1, 2, 3, 3, 1])
numbers.count_in_bins(how_many_bins=3)
\end{lstlisting}

\unsection{Yes or no, and conditions}{from-numpy:conditions}

\textbf{The job.} Ask a question of every number and use the answers: count them, keep the numbers
where the answer is yes, or replace them (NumPy's comparisons, boolean masks, \code{where},
\code{count\_nonzero}, \code{any} and \code{all}). In the guide:
\href{https://numpy.org/doc/stable/user/basics.indexing.html}{Indexing on ndarrays}, ``Boolean array
indexing''.

\textbf{In simplydata.} \code{turn\_into\_yes\_or\_no(is\_above=)} is the mask;
\code{keep\_numbers\_where}, \code{replace\_numbers\_where} and \code{count\_numbers\_that} use it, and
\code{check\_if\_any\_number} and \code{check\_if\_all\_numbers} are \code{any} and \code{all}.

\begin{lstlisting}[style=py]
rain = sdata.number_list([0.0, 2.5, 0.0, 12.1, 4.4, 0.0])
rain.count_numbers_that(is_above=0)      # 3 days with rain
wet = rain.turn_into_yes_or_no(is_above=0)
rain.check_if_any_number(is_above=10)
rain.replace_numbers_where(wet, with_value=1)
print(rain)
\end{lstlisting}

\unsection{Empty values}{from-numpy:empty}

\textbf{The job.} Find the missing values (\code{nan}), work round them, or fill them (NumPy's
\code{isnan}, the \code{nan} functions such as \code{nanmean}, and \code{nan\_to\_num}). In the guide:
\href{https://numpy.org/doc/stable/user/absolute_beginners.html}{the absolute beginners guide},
``More useful array operations''.

\textbf{In simplydata.} An empty value is written \code{None} or \code{float("nan")}.
\code{count\_empty\_numbers} counts them, the summaries leave them out and say so, and
\code{fill\_empty\_numbers} fills them.

\begin{lstlisting}[style=py]
readings = sdata.number_list([4.1, None, 3.9, 4.4, None])
readings.count_empty_numbers()
readings.find_average()                  # of the 3 numbers that are there
readings.fill_empty_numbers(with_value=0)
\end{lstlisting}

\unsection{Random numbers}{from-numpy:random}

\textbf{The job.} Make random numbers from a distribution, pick at random, and get the same numbers
again with a seed (NumPy's \code{default\_rng} and its Generator methods). In the guides:
\href{https://numpy.org/doc/stable/user/absolute_beginners.html}{the absolute beginners guide},
``Generating random numbers'', and \href{https://numpy.org/doc/stable/reference/random/index.html}%
{Random sampling}.

\textbf{In simplydata.} \code{sdata.make\_random\_numbers} makes normal or uniform numbers (and whole
numbers); \code{sdata.distribution("poisson", average=3).make\_random\_numbers(how\_many=10)} makes them
from any of 28 distributions. \code{numbers.make\_random\_picks} is \code{choice}. \code{with\_seed=}
is the seed.

\begin{lstlisting}[style=py]
heights = sdata.make_random_numbers(how_many=100, average=150, standard_deviation=8, with_seed=1)
dice = sdata.make_random_numbers(how_many=10, between=[1, 6], whole_numbers=True, with_seed=1)
arrivals = sdata.distribution("poisson", average=3).make_random_numbers(how_many=10, with_seed=1)
picked = heights.make_random_picks(how_many=5, with_seed=1)
print(picked)
\end{lstlisting}

\unsection{Saving and reading numbers}{from-numpy:files}

\textbf{The job.} Write numbers to a file and read them back (NumPy's \code{save}, \code{savetxt},
\code{loadtxt} and \code{genfromtxt}). In the guide: \href{https://numpy.org/doc/stable/user/basics.io.html}%
{I/O with NumPy}.

\textbf{In simplydata.} \code{numbers.save("file.csv")} writes CSV, Excel, Parquet, JSON or a sound
file, as the end of the name says. \code{sdata.number\_list("file.csv")} reads them back, and
\code{empty\_values=[-999]} is \code{genfromtxt}'s \code{missing\_values}. A netCDF file opens as a
block with \code{sdata.number\_block("file.nc")}.

\begin{lstlisting}[style=py]
numbers = sdata.number_list([4.1, 3.9, -999, 4.4])
numbers.save("readings.csv")
again = sdata.number_list("readings.csv", empty_values=[-999])
again.count_empty_numbers()
\end{lstlisting}

\unsection{Grids as matrices}{from-numpy:matrices}

\textbf{The job.} The determinant, the inverse and the eigenvalues of a square grid, and solving the
equations grid times x = b (NumPy's \code{np.linalg}). In the guides:
\href{https://numpy.org/doc/stable/user/absolute_beginners.html}{the absolute beginners guide},
``Working with mathematical formulas'', and
\href{https://numpy.org/doc/stable/reference/routines.linalg.html}{Linear algebra}.

\textbf{In simplydata.} \code{find\_determinant}, \code{take\_inverse} and \code{find\_eigenvalues};
\code{solve\_for(b)} is \code{solve}, and on a grid that is not square it gives the least-squares answer
(\code{lstsq}) and says so. \code{exactly=True} on each gives SymPy's exact answer, as fractions. To
solve equations, \code{solve\_for} is better than an inverse. \code{multiply\_as\_matrix\_by} is the
matrix product (\code{@}, \code{np.matmul}); \code{find\_rank}, \code{find\_size} and
\code{find\_condition\_number} are \code{matrix\_rank}, \code{norm} and \code{cond}.

\begin{lstlisting}[style=py]
grid = sdata.number_grid([[2, 1], [1, 3]])
grid.find_determinant()                  # 5
grid.find_eigenvalues()
x = grid.solve_for([5, 10])              # x = [1, 3]
print(x)
grid.solve_for([1, 1], exactly=True)     # [2/5, 1/5]
grid.multiply_as_matrix_by(sdata.number_grid([[0, 1], [1, 0]]))
grid.find_rank()                         # 2
\end{lstlisting}

\unsection{Fourier transforms}{from-numpy:fourier}

\textbf{The job.} Turn readings in time into frequencies, and back (NumPy's and SciPy's \code{fft},
\code{rfft}, \code{fftfreq} and \code{ifft}). In the guide:
\href{https://docs.scipy.org/doc/scipy/tutorial/fft.html}{Discrete Fourier Transforms (scipy.fft)}.

\textbf{In simplydata.} \code{numbers.take\_fourier\_transform(time\_step=)} gives a table with each
frequency in hertz, its size and its phase: \code{rfft} and \code{rfftfreq} in one word.
\code{sdata.find\_inverse\_fourier\_transform(table)} goes back. \code{take\_cosine\_transform} is
\code{dct}. \code{find\_spectrum} gives the power at each frequency.

\begin{lstlisting}[style=py]
wave = sdata.make_wave("sine", frequency=5, time_step=0.01, how_many=200)
frequencies = wave.take_fourier_transform(time_step=0.01)
back = sdata.find_inverse_fourier_transform(frequencies)
back.count()
\end{lstlisting}

\unsection{Areas, volumes and paths over time}{from-numpy:integration}

\textbf{The job.} The area under a curve, the volume under a surface, the running area under readings,
and how a quantity changes over time from a rate (SciPy's \code{quad}, \code{dblquad},
\code{cumulative\_trapezoid} and \code{solve\_ivp}). In the guide:
\href{https://docs.scipy.org/doc/scipy/tutorial/integrate.html}{Integration (scipy.integrate)}.

\textbf{In simplydata.} \code{sdata.find\_area\_under}, \code{sdata.find\_volume\_under},
\code{numbers.turn\_into\_running\_area} and \code{sdata.find\_path\_over\_time}. The formula is written
as text, with letters. \code{sdata.find\_path\_between\_ends} is \code{solve\_bvp}, for a path with a value
given at each end, and \code{sdata.find\_area\_points} gives the points and weights of Gaussian quadrature
(\code{roots\_legendre} and its kin).

\begin{lstlisting}[style=py]
sdata.find_area_under("exp(-x**2)", between=[0, "infinity"])
sdata.find_volume_under("x * y", between={"x": [0, 1], "y": [0, 2]})
flow = sdata.number_list([0.0, 1.5, 2.0, 1.0, 0.5])
flow.turn_into_running_area(time_step=60)
path = sdata.find_path_over_time("dN/dt = r * N", at_start={"N": 10}, between=[0, 10], known={"r": 0.3})
\end{lstlisting}

\unsection{Smallest, biggest and solutions}{from-numpy:optimization}

\textbf{The job.} Find where a formula is smallest or biggest, solve equations, fit a curve to data,
and solve least squares with limits (SciPy's \code{minimize}, \code{root}, \code{curve\_fit},
\code{linprog}, \code{nnls} and \code{lsq\_linear}). In the guide:
\href{https://docs.scipy.org/doc/scipy/tutorial/optimize.html}{Optimization (scipy.optimize)}.

\textbf{In simplydata.} \code{sdata.find\_where\_smallest} and \code{sdata.find\_where\_biggest} (with
\code{rules=} for linear programming), \code{sdata.solve}, \code{table.fit\_formula} for curves, and
\code{grid.solve\_for(b, between=[0, None])} for answers that may not be negative, as in a diet mix or
the unmixing of spectra. \code{sdata.find\_best\_pairing(costs=grid)} is \code{linear\_sum\_assignment}: it
pairs each row with one column so that the total cost is smallest.

\begin{lstlisting}[style=py]
sdata.find_where_smallest("x**2 - 4*x + 1", between=[0, 5])
sdata.find_where_smallest("2*x + 3*y", rules=["x + y >= 4", "x >= 0", "y >= 0"])
sdata.solve("x**3 - 2*x = 5", between=[0, 3])
# one row per nutrient (protein, fibre, fat), one column per food: grams of the nutrient in one portion
foods = sdata.number_grid([[20, 2, 5], [10, 15, 1], [2, 8, 30]])
amounts = foods.solve_for([60, 40, 50], between=[0, None])          # portions of each food, never below 0
grid = sdata.number_grid([[4, 1, 3], [2, 0, 5], [3, 2, 2]])         # cost of each worker (row) for each job
sdata.find_best_pairing(costs=grid)
\end{lstlisting}

\unsection{Values between the points}{from-numpy:interpolation}

\textbf{The job.} Find values between measured points, on a line, on a map or in a block (SciPy's
\code{interp1d}, \code{CubicSpline}, \code{griddata} and \code{RegularGridInterpolator}). In the guide:
\href{https://docs.scipy.org/doc/scipy/tutorial/interpolate.html}{Interpolation (scipy.interpolate)}.

\textbf{In simplydata.} \code{table.find\_values\_at} on a table's columns,
\code{table.make\_map\_of} for scattered points on a map, and \code{block.find\_values\_at} inside a
labelled block. \code{sdata.curve\_through\_points(x=..., y=...)} is a spline you can keep (SciPy's
\code{make\_interp\_spline} and \code{CubicSpline}): read it anywhere with \code{curve.find\_value\_at}, and
take its slope, area and zeros with \code{find\_slope\_at}, \code{find\_area\_between} and \code{find\_zeros}.

\begin{lstlisting}[style=py]
growth = sdata.get_example("growth")
growth.keep_rows_where("country", is_equal_to="Applia")
growth.find_values_at("population", at={"year": [2001.5, 2003.5]})
curve = sdata.curve_through_points(x=[0, 5, 10, 20], y=[8.1, 7.9, 6.2, 5.0])
curve.find_value_at(12)
curve.find_slope_at(12)
\end{lstlisting}

\unsection{Signals}{from-numpy:signals}

\textbf{The job.} Smooth readings, keep or remove frequencies with a filter, find peaks, and find the
spectrum (SciPy's \code{savgol\_filter}, \code{butter} with \code{filtfilt}, \code{find\_peaks} and
\code{welch}). In the guide: \href{https://docs.scipy.org/doc/scipy/tutorial/signal.html}{Signal
Processing (scipy.signal)}.

\textbf{In simplydata.} \code{smooth}, \code{keep\_each\_frequency\_that} and
\code{remove\_each\_frequency\_that}, \code{find\_each\_peak\_that} and \code{find\_spectrum}. Each
takes \code{time\_step=}, the time between readings, so frequencies are in hertz. The filter designs of
\code{scipy.signal} (Butterworth, Chebyshev, elliptic, Bessel and FIR filters) are values of
\code{sdata.frequency\_filter}; \code{system.move\_poles\_to} is \code{place\_poles}, and
\code{find\_phase\_locking\_to} measures how steady the phase is.

\begin{lstlisting}[style=py]
signal = sdata.make_wave("sine", frequency=2, time_step=0.01, how_many=300)
noise = sdata.make_random_numbers(how_many=300, average=0, standard_deviation=0.3, with_seed=1)
signal.add(noise)
signal.keep_each_frequency_that(is_below=5, time_step=0.01)
signal.find_each_peak_that(is_at_least=0.5)
\end{lstlisting}

\unsection{More linear algebra}{from-numpy:linalg}

\textbf{The job.} Break a matrix down into parts (SVD, QR, LU, Cholesky, Schur and more), the matrix
exponential, and the null space (SciPy's \code{scipy.linalg}). In the guide:
\href{https://docs.scipy.org/doc/scipy/tutorial/linalg.html}{Linear Algebra (scipy.linalg)}.

\textbf{In simplydata.} \code{grid.break\_down(test\_method="svd")} gives the parts as a dict of
grids and prints what each part is; the other methods are \code{"qr"}, \code{"lu"},
\code{"cholesky"}, \code{"schur"}, \code{"polar"}, \code{"ldl"}, \code{"hessenberg"} and \code{"qz"}.
\code{take\_matrix\_exp} is \code{expm} and \code{find\_null\_space} is \code{null\_space}. The other
matrix functions are one family word, \code{grid.take\_matrix\_function("log")} (\code{logm}, \code{sqrtm},
\code{cosm} ...); \code{check\_if\_grid\_is("symmetric")} checks a kind of grid, and
\code{sdata.solve\_matrix\_equation} solves the Sylvester, Lyapunov and Riccati equations.

\begin{lstlisting}[style=py]
grid = sdata.number_grid([[4, 2], [2, 3]])
parts = grid.break_down(test_method="cholesky")
print(parts["U"])
parts = grid.break_down(test_method="svd")
print(parts["singular values"])
grid.find_null_space()
grid.take_matrix_function("square root")                 # SciPy's sqrtm
grid.check_if_grid_is("symmetric")
\end{lstlisting}

\unsection{Points in space}{from-numpy:spatial}

\textbf{The job.} The nearest points, distances between points, and the fence round a set of points
(SciPy's \code{KDTree}, \code{distance.cdist} and \code{ConvexHull}). In the guide:
\href{https://docs.scipy.org/doc/scipy/tutorial/spatial.html}{Spatial data structures and algorithms
(scipy.spatial)}.

\textbf{In simplydata.} These are table words, because points are rows with an east and a north:
\code{table.find\_nearest\_rows\_to}, \code{table.find\_distances\_to\_rows\_of} and
\code{table.find\_fence\_around}. \code{table.find\_triangles\_between\_points} is \code{Delaunay} and
\code{table.find\_regions\_around\_points} is \code{Voronoi}, with the area of each region.

\begin{lstlisting}[style=py]
island = sdata.get_example("island")
island.find_nearest_rows_to({"east": 4, "north": 1}, how_many=3)
island.find_fence_around(x="east", y="north")
island = sdata.get_example("island")
island.find_triangles_between_points(x="east", y="north")
\end{lstlisting}

\unsection{Special functions}{from-numpy:special}

\textbf{The job.} The gamma and beta functions, the error function, Bessel functions, and the number of
ways to choose (SciPy's \code{scipy.special}). In the guide:
\href{https://docs.scipy.org/doc/scipy/tutorial/special.html}{Special Functions (scipy.special)}.

\textbf{In simplydata.} A formula can use \code{gamma}, \code{lgamma}, \code{digamma}, \code{beta},
\code{erf}, \code{erfc} and \code{besselJ}. \code{sdata.count\_ways\_to\_choose} is \code{comb} and
\code{sdata.count\_ways\_to\_order} is \code{perm}; \code{take\_log\_odds} and \code{turn\_into\_chances}
are \code{logit} and \code{expit}. Every other function of \code{scipy.special} is a name of
\code{numbers.take\_function} (see ``Maths on every number''), and the same names work in a formula with
\code{\_} for spaces, such as \code{bessel\_j(1, x)}. \code{sdata.find\_zeros\_of("bessel J")} gives the zeros
(\code{jn\_zeros} and its kin), and \code{sdata.polynomial("legendre", degree=3)} makes the polynomials of
\code{scipy.special} and \code{np.polynomial} as formulas.

\begin{lstlisting}[style=py]
formula = sdata.formula("gamma(x)")
formula.find_value_at(x=4.5)
sdata.count_ways_to_choose(3, out_of=10)                 # 120
chances = sdata.number_list([0.1, 0.5, 0.9])
chances.take_log_odds()
sdata.find_zeros_of("bessel J", how_many=3)
formula = sdata.polynomial("legendre", degree=3)
formula.find_value_at(x=0.5)
\end{lstlisting}

\unsection{Statistics}{from-numpy:statistics}

\textbf{The job.} Use a distribution (density, chances, percentiles, random numbers), fit one to data,
and run tests (SciPy's \code{scipy.stats}). In the guide:
\href{https://docs.scipy.org/doc/scipy/tutorial/stats.html}{Statistics (scipy.stats)}.

\textbf{In simplydata.} \code{sdata.distribution(...)} with \code{find\_density\_at} (\code{pdf}),
\code{find\_chance\_that\_value} (\code{cdf} and \code{sf}) and \code{find\_percentile} (\code{ppf});
\code{numbers.fit\_distribution} is \code{fit}. The tests are table words, such as
\code{table.compare\_averages\_of}; the rest of this manual explains them. Every distribution of
\code{scipy.stats} has a plain name, such as \code{sdata.distribution("skew normal", shape=4)}.

\begin{lstlisting}[style=py]
distribution = sdata.distribution("normal", average=50, standard_deviation=10)
distribution.find_chance_that_value(is_below=60)
distribution.find_percentile(95)
waits = sdata.make_random_numbers(how_many=200, average=10, standard_deviation=2, with_seed=1)
waits.fit_distribution("normal")
\end{lstlisting}

\unsection{Images and grids of cells}{from-numpy:ndimage}

\textbf{The job.} Smooth a grid, find patches of cells that pass a test, turn and resize a grid
(SciPy's \code{scipy.ndimage}: \code{gaussian\_filter}, \code{label}, \code{rotate} and \code{zoom}).
In the guide: \href{https://docs.scipy.org/doc/scipy/tutorial/ndimage.html}{Multidimensional image
processing (scipy.ndimage)}.

\textbf{In simplydata.} \code{grid.smooth}, \code{grid.find\_each\_patch\_that}, \code{grid.rotate} and
\code{grid.change\_size\_to}; on a block, \code{block.smooth(over=[...])}. \code{grid.find\_edges()} is
\code{sobel}; \code{grid.change\_patch\_shapes("grow")} is one family word for \code{binary\_dilation},
\code{binary\_erosion}, opening, closing and filling holes; \code{find\_distance\_to\_nearest} is
\code{distance\_transform\_edt}.

\begin{lstlisting}[style=py]
grid = sdata.number_grid([[0, 0, 1, 1], [0, 0, 1, 1], [5, 0, 0, 0], [5, 5, 0, 0]])
grid.find_each_patch_that(is_above=0)
grid.change_size_to(rows=8, columns=8)
grid.rotate(by_degrees=15)
grid = sdata.number_grid([[0, 0, 0, 0], [0, 1, 1, 0], [0, 1, 1, 0], [0, 0, 0, 0]])
grid.find_edges()
\end{lstlisting}

\unsection{Sound files and other files}{from-numpy:file-io}

\textbf{The job.} Read and write sound files, and other file types (SciPy's \code{scipy.io}, such as
\code{wavfile}). In the guide: \href{https://docs.scipy.org/doc/scipy/tutorial/io.html}{File IO
(scipy.io)}.

\textbf{In simplydata.} \code{numbers.save("sound.wav", time\_step=)} and
\code{sdata.number\_list("sound.wav")}. netCDF files open as a block, and
\code{sdata.show\_whats\_inside("file.nc")} lists what is in one.

\begin{lstlisting}[style=py]
tone = sdata.make_wave("sine", frequency=440, time_step=1 / 8000, how_many=8000)
tone.times(0.5)
tone.save("tone.wav", time_step=1 / 8000)
again = sdata.number_list("tone.wav")
again.count()
\end{lstlisting}

\unsection{The jobs left to NumPy and SciPy}{from-numpy:left-out}

simplydata does the jobs most people need in research and teaching, with one word each. Many jobs that
were once left to NumPy now have words: \code{np.diff} is \code{turn\_into\_changes}, \code{np.histogram} is
\code{count\_in\_bins}, \code{np.argsort} is \code{sort(show="positions")}, \code{np.meshgrid} is
\code{sdata.make\_coordinate\_grids}, \code{np.linalg.norm} is \code{find\_size}, the matrix product is
\code{multiply\_as\_matrix\_by}, and \code{np.polynomial} is \code{sdata.polynomial}. These jobs are left out
on purpose. Each one is a hand-over: \code{array = numbers.as\_numpy()} (or \code{grid.as\_numpy()},
\code{block.as\_numpy()}), then NumPy or SciPy code, then \code{sdata.number\_list(...)},
\code{sdata.number\_grid(...)} or \code{sdata.number\_block(...)} to come back.
\begin{itemize}
  \item \textbf{Sums over named indices} (\code{np.einsum}, \code{np.linalg.tensorinv}): einsum is its own
    notation; \code{multiply\_as\_matrix\_by} covers the usual uses.
  \item \textbf{Data types, copies and views, structured arrays and strings in arrays}: simplydata keeps
    numbers as decimals and makes copies, so these questions do not come up.
  \item \textbf{Sparse arrays} (\code{scipy.sparse}) and their eigenvalue and graph routines.
  \item \textbf{The intersection of half-spaces} (\code{scipy.spatial.HalfspaceIntersection}), expert
    geometry with no data job.
  \item \textbf{Steps inside a solver} (\code{qr\_update}, \code{expm\_frechet}, the BLAS and LAPACK
    functions): \code{sdata.help} names the word that runs them.
\end{itemize}

One more example: a matrix product, and back.

\begin{lstlisting}[style=py]
a = sdata.number_grid([[1, 2], [3, 4]])
b = sdata.number_grid([[0, 1], [1, 0]])
product = a.as_numpy() @ b.as_numpy()    # NumPy's matrix product
grid = sdata.number_grid(product)        # back to simplydata
grid.find_determinant()
\end{lstlisting}
).
% The topics follow the order of the NumPy absolute beginners guide and the NumPy user guide (fundamentals), then the
% SciPy user guide, used as a checklist only: the text and the examples are our own, and each section links to the
% official page. Every example in a py box is run by tests/test_from_numpy_chapter.py (with "import simplydata as
% sdata" in front), so it cannot go out of date.

\unchapter{From NumPy and SciPy to simplydata}{from-numpy}

This chapter is for you if you already know NumPy or SciPy, or if you want to learn them after
simplydata. It follows the topics of the official guides, in their order: the
\href{https://numpy.org/doc/stable/user/absolute_beginners.html}{NumPy absolute beginners guide} and
the \href{https://numpy.org/doc/stable/user/basics.html}{NumPy fundamentals}, then the
\href{https://docs.scipy.org/doc/scipy/tutorial/index.html}{SciPy user guide}. For each topic it says
what the job is, which simplydata words do it, and a short example. Some jobs are left to NumPy and
SciPy on purpose. For those, the chapter shows how to hand the numbers over and how to bring the
answer back.

Every example starts with \code{import simplydata as sdata}. It is written once here, not in each box.

\unsection{Five things to know first}{from-numpy:first}

\begin{enumerate}
  \item \textbf{Three makers, one family.} \code{sdata.number\_list} makes a list of numbers (NumPy's
    1-D array), \code{sdata.number\_grid} makes rows and columns (a 2-D array) and
    \code{sdata.number\_block} makes numbers in any number of directions, each with a name (an n-D
    array, or an xarray DataArray). The words of numbers work on grids and blocks too.
  \item \textbf{Positions count from 1.} \code{numbers.get\_number\_at(1)} is the first number. NumPy
    counts from 0, so the sentence a word prints gives NumPy's position too.
  \item \textbf{A word changes the object, and \code{undo} takes it back.} \code{numbers.take\_log()}
    changes \code{numbers} itself and prints what it did. \code{numbers.undo()} brings back the numbers
    from before. A word that finds something (an average, a determinant) gives it back and changes
    nothing.
  \item \textbf{Empty values stay empty.} NumPy writes a missing value as \code{nan}. simplydata calls it
    an empty value, says how many there are, and leaves them out of averages and totals, as NumPy's
    \code{nanmean} does.
  \item \textbf{The door is always open.} \code{array = numbers.as\_numpy()} gives a plain NumPy array,
    for any NumPy or SciPy call. \code{sdata.number\_list(array)}, \code{sdata.number\_grid(array)} and
    \code{sdata.number\_block(array)} bring the answer back.
\end{enumerate}

\begin{lstlisting}[style=py]
numbers = sdata.number_list([2.5, 3.1, 4.8, 5.0, 6.2])
array = numbers.as_numpy()               # a NumPy array: use any NumPy or SciPy code here
import numpy as np
steps = np.diff(array)
back = sdata.number_list(steps)          # and back to simplydata
back.find_average()
\end{lstlisting}

Do you know a NumPy, SciPy or SymPy name, but not the simplydata word? Ask:
\code{sdata.help("np.mean")}, \code{sdata.help("np.linalg.solve")}, \code{sdata.help("rng.normal")}
or \code{sdata.help("scipy.stats.gamma")}. The name can have \code{np.}, \code{np.linalg.},
\code{scipy.}, \code{stats.} or \code{optimize.} in front, and the answer is the word for that package:
\code{np.mean} is the average of numbers, but a bare \code{mean} is the average of a table column.

\unsection{Making numbers, grids and blocks}{from-numpy:making}

\textbf{The job.} Make an array: from a list, as a run of evenly spaced numbers, filled with one
value, or of a given shape. In the guides:
\href{https://numpy.org/doc/stable/user/absolute_beginners.html}{the absolute beginners guide}, ``How
to create a basic array'', and \href{https://numpy.org/doc/stable/user/basics.creation.html}{Array
creation}.

\textbf{In simplydata.} \code{sdata.number\_list([...])} is \code{np.array}; \code{sdata.number\_list(0,
to=1, how\_many=11)} is \code{np.linspace} and \code{step=} is \code{np.arange}.
\code{sdata.number\_grid(rows=2, columns=3, filled\_with=0)} is \code{np.zeros} (and \code{np.full});
\code{sdata.number\_block(of\_shape=(2, 3, 4), filled\_with=1)} is \code{np.ones} with any shape.
Special grids come from \code{kind=}: \code{sdata.number\_grid(kind="identity", size=3)} is \code{np.eye},
\code{kind="diagonal"} is \code{np.diag}, and \code{"toeplitz"}, \code{"hilbert"} or \code{"pascal"} are the
grids of \code{scipy.linalg}. \code{sdata.make\_coordinate\_grids(x=..., y=...)} is \code{np.meshgrid}.

\begin{lstlisting}[style=py]
temperatures = sdata.number_list([8.3, 10.3, 11.0, 9.2, 7.5])
steps = sdata.number_list(0, to=1, how_many=11)          # 0, 0.1, ..., 1
every_five = sdata.number_list(0, to=50, step=5)
grid = sdata.number_grid([[1, 2, 3], [4, 5, 6]])
zeros = sdata.number_grid(rows=2, columns=3, filled_with=0)
block = sdata.number_block(of_shape=(2, 3, 4), filled_with=1)
print(grid)
grid = sdata.number_grid(kind="identity", size=3)       # NumPy's eye
a, b = sdata.make_coordinate_grids(x=[0, 1, 2], y=[10, 20])    # NumPy's meshgrid
\end{lstlisting}

\unsection{Shape and size}{from-numpy:shape}

\textbf{The job.} Know how many numbers there are, how many rows and columns, and how many directions
(NumPy's \code{size}, \code{shape} and \code{ndim}). In the guide:
\href{https://numpy.org/doc/stable/user/absolute_beginners.html}{the absolute beginners guide}, ``How
do you know the shape and size of an array?''.

\textbf{In simplydata.} \code{numbers.count()} is \code{size}; \code{grid.find\_shape()} and
\code{block.find\_shape()} are \code{shape}; \code{block.count\_directions()} is \code{ndim}.
\code{grid.count\_rows()} and \code{grid.count\_columns()} say one size each.

\begin{lstlisting}[style=py]
grid = sdata.number_grid([[1, 2, 3], [4, 5, 6]])
grid.find_shape()                        # (2, 3)
grid.count_rows()
block = sdata.number_block(of_shape=(2, 3, 4), filled_with=0)
block.count_directions()                 # 3
\end{lstlisting}

\unsection{Picking numbers out}{from-numpy:indexing}

\textbf{The job.} Take one number, a slice, a row or a column, or the numbers that pass a test
(NumPy's indexing and slicing, and boolean indexing). In the guides:
\href{https://numpy.org/doc/stable/user/absolute_beginners.html}{the absolute beginners guide},
``Indexing and slicing'', and \href{https://numpy.org/doc/stable/user/basics.indexing.html}{Indexing
on ndarrays}.

\textbf{In simplydata.} \code{get\_number\_at} is \code{a[i]} (from 1); \code{keep\_numbers\_by\_position(2,
to=4)} is a slice and keeps both ends; \code{grid.get\_row} and \code{grid.get\_column} are \code{a[i, :]}
and \code{a[:, j]}. \code{keep\_each\_number\_that(is\_above=5)} is \code{a[a > 5]}.
\code{grid.get\_part("diagonal")} and \code{grid.get\_part("upper triangle")} are \code{np.diagonal} and
\code{np.triu}; \code{find\_positions\_where(is\_above=5)} is \code{np.nonzero} (counting from 1).

\begin{lstlisting}[style=py]
numbers = sdata.number_list([4, 8, 15, 16, 23, 42])
numbers.get_number_at(1)                 # 4, the first number (NumPy's a[0])
numbers.keep_numbers_by_position(2, to=4)
print(numbers)                           # 8, 15, 16
numbers.undo()
numbers.keep_each_number_that(is_above=10)
grid = sdata.number_grid([[1, 2, 3], [4, 5, 6]])
grid.get_row(2)
grid.get_column(3)
grid.get_part("diagonal")                # 1 and 5
\end{lstlisting}

\unsection{Maths on every number}{from-numpy:maths}

\textbf{The job.} Add, multiply, take logs and square roots of every number at once (NumPy's universal
functions), also between two arrays of the same length. In the guides:
\href{https://numpy.org/doc/stable/user/absolute_beginners.html}{the absolute beginners guide}, ``Basic
array operations'', and \href{https://numpy.org/doc/stable/user/basics.ufuncs.html}{Universal
functions}.

\textbf{In simplydata.} \code{add}, \code{subtract}, \code{times}, \code{divide\_by},
\code{raise\_to\_power}, \code{take\_square\_root}, \code{take\_log}, \code{take\_exp},
\code{take\_sin}, \code{round\_to} and their family. Each takes one number or numbers of the same
length. \code{+}, \code{-}, \code{*} and \code{/} work too and give new numbers.
Every other function of NumPy and \code{scipy.special} is one family word:
\code{numbers.take\_function("hyperbolic sine")} is \code{np.sinh}, and
\code{take\_function("bessel J", order=1)} is \code{special.jv}. The library names work too, and
\code{numbers.take\_function()} with no name lists them all. \code{work\_on\_bits} does the bit operations
(\code{np.bitwise\_and} and the rest). \code{turn\_into\_changes} is \code{np.diff},
\code{turn\_into\_slopes} is \code{np.gradient} and \code{find\_product} is \code{np.prod}.

\begin{lstlisting}[style=py]
prices = sdata.number_list([1.20, 3.50, 2.25])
prices.times(1.25)                       # every price with 25 % added
prices.round_to(with_decimals=2)
counts = sdata.number_list([10, 100, 1000])
counts.take_log(base=10)
both = sdata.number_list([1, 2, 3]) + sdata.number_list([10, 20, 30])
print(both)
numbers = sdata.number_list([0.5, 1.0, 2.0])
numbers.take_function("hyperbolic sine")                 # NumPy's sinh
numbers.undo()
numbers.take_function("bessel J", order=1)               # SciPy's special.jv
numbers = sdata.number_list([3, 5, 9, 10])
numbers.turn_into_changes()                              # 2, 4, 1: NumPy's diff
\end{lstlisting}

\unsection{Stretching to the same shape}{from-numpy:broadcasting}

\textbf{The job.} Work on two arrays of different shapes by stretching the smaller one (NumPy's
broadcasting). In the guide: \href{https://numpy.org/doc/stable/user/basics.broadcasting.html}%
{Broadcasting}.

\textbf{In simplydata.} The maths words stretch by themselves, as NumPy does, and print a note when
they do. \code{block.show\_how\_it\_lines\_up\_with(other)} explains how two blocks line up before you
use them, and \code{sdata.stretch\_to\_same\_shape(a, b)} gives both at the same shape.
\code{sdata.hide\_stretching\_notes()} turns the notes off.

\begin{lstlisting}[style=py]
a = sdata.number_block(of_shape=(3, 4), filled_with=1)
b = sdata.number_block(of_shape=(1, 4), filled_with=2)
a.show_how_it_lines_up_with(b)
a_big, b_big = sdata.stretch_to_same_shape(a, b)
b_big.find_shape()
\end{lstlisting}

\unsection{Summaries}{from-numpy:summaries}

\textbf{The job.} The total, average, smallest and biggest number, the standard deviation and
percentiles, of all numbers or along one direction (NumPy's aggregation functions with
\code{axis=}). In the guide: \href{https://numpy.org/doc/stable/user/absolute_beginners.html}{the
absolute beginners guide}, ``More useful array operations''.

\textbf{In simplydata.} \code{find\_total}, \code{find\_average}, \code{find\_median},
\code{find\_smallest}, \code{find\_biggest}, \code{find\_standard\_deviation} and
\code{find\_percentile}. On a block, \code{find\_average(over="month")} averages along a named
direction, as \code{axis=} does with a number. The standard deviation is the sample one (NumPy's
\code{ddof=1}), as in statistics books.

\begin{lstlisting}[style=py]
numbers = sdata.number_list([8.3, 10.3, 11.0, 9.2, 7.5])
numbers.find_average()
numbers.find_standard_deviation()        # sample standard deviation (NumPy's ddof=1)
numbers.find_percentile(90)
block = sdata.number_block(of_shape=(2, 3), filled_with=1, directions_named=["year", "month"])
block.find_average(over="month")
\end{lstlisting}

\unsection{Changing the shape}{from-numpy:reshaping}

\textbf{The job.} Reshape, turn rows into columns, join arrays and split them, add or remove a
direction (NumPy's \code{reshape}, \code{transpose}, \code{concatenate}, \code{split},
\code{expand\_dims} and \code{squeeze}). In the guide:
\href{https://numpy.org/doc/stable/user/absolute_beginners.html}{the absolute beginners guide}, ``Can
you reshape an array?'' and ``Transposing and reshaping a matrix''.

\textbf{In simplydata.} \code{numbers.make\_number\_grid(rows=, columns=)} makes a grid from numbers;
\code{block.change\_shape\_to(...)} changes a block's own shape. \code{grid.swap\_rows\_and\_columns()}
is \code{.T}. \code{block.add\_numbers\_from(other, over=)} joins, \code{block.split\_into\_pieces}
splits, and \code{add\_direction} and \code{remove\_directions\_of\_size\_one} are \code{expand\_dims}
and \code{squeeze}.

\begin{lstlisting}[style=py]
numbers = sdata.number_list([1, 2, 3, 4, 5, 6])
grid = numbers.make_number_grid(rows=2, columns=3)
grid.swap_rows_and_columns()
grid.find_shape()                        # (3, 2)
block = sdata.number_block(of_shape=(2, 6), filled_with=0, directions_named=["row", "column"])
block.change_shape_to(2, 3, 2)
block.find_shape()
\end{lstlisting}

\unsection{Sorting, unique numbers and counts}{from-numpy:sorting}

\textbf{The job.} Sort, keep each different number once, and count how often each comes (NumPy's
\code{sort}, \code{unique} and \code{unique\_counts}). In the guide:
\href{https://numpy.org/doc/stable/user/absolute_beginners.html}{the absolute beginners guide},
``Adding, removing, and sorting elements'' and ``How to get unique items and counts''.

\textbf{In simplydata.} \code{sort} (\code{biggest\_first=True} for the other way),
\code{remove\_duplicate\_numbers} and \code{count\_each}, which gives a table of each number and its
count. \code{sort(show="positions")} is \code{np.argsort}, and \code{count\_in\_bins(how\_many\_bins=3)} is
\code{np.histogram} as numbers.

\begin{lstlisting}[style=py]
numbers = sdata.number_list([3, 1, 2, 3, 3, 1])
numbers.count_each()
numbers.sort(biggest_first=True)
numbers.remove_duplicate_numbers()
print(numbers)
numbers = sdata.number_list([3, 1, 2, 3, 3, 1])
numbers.count_in_bins(how_many_bins=3)
\end{lstlisting}

\unsection{Yes or no, and conditions}{from-numpy:conditions}

\textbf{The job.} Ask a question of every number and use the answers: count them, keep the numbers
where the answer is yes, or replace them (NumPy's comparisons, boolean masks, \code{where},
\code{count\_nonzero}, \code{any} and \code{all}). In the guide:
\href{https://numpy.org/doc/stable/user/basics.indexing.html}{Indexing on ndarrays}, ``Boolean array
indexing''.

\textbf{In simplydata.} \code{turn\_into\_yes\_or\_no(is\_above=)} is the mask;
\code{keep\_numbers\_where}, \code{replace\_numbers\_where} and \code{count\_numbers\_that} use it, and
\code{check\_if\_any\_number} and \code{check\_if\_all\_numbers} are \code{any} and \code{all}.

\begin{lstlisting}[style=py]
rain = sdata.number_list([0.0, 2.5, 0.0, 12.1, 4.4, 0.0])
rain.count_numbers_that(is_above=0)      # 3 days with rain
wet = rain.turn_into_yes_or_no(is_above=0)
rain.check_if_any_number(is_above=10)
rain.replace_numbers_where(wet, with_value=1)
print(rain)
\end{lstlisting}

\unsection{Empty values}{from-numpy:empty}

\textbf{The job.} Find the missing values (\code{nan}), work round them, or fill them (NumPy's
\code{isnan}, the \code{nan} functions such as \code{nanmean}, and \code{nan\_to\_num}). In the guide:
\href{https://numpy.org/doc/stable/user/absolute_beginners.html}{the absolute beginners guide},
``More useful array operations''.

\textbf{In simplydata.} An empty value is written \code{None} or \code{float("nan")}.
\code{count\_empty\_numbers} counts them, the summaries leave them out and say so, and
\code{fill\_empty\_numbers} fills them.

\begin{lstlisting}[style=py]
readings = sdata.number_list([4.1, None, 3.9, 4.4, None])
readings.count_empty_numbers()
readings.find_average()                  # of the 3 numbers that are there
readings.fill_empty_numbers(with_value=0)
\end{lstlisting}

\unsection{Random numbers}{from-numpy:random}

\textbf{The job.} Make random numbers from a distribution, pick at random, and get the same numbers
again with a seed (NumPy's \code{default\_rng} and its Generator methods). In the guides:
\href{https://numpy.org/doc/stable/user/absolute_beginners.html}{the absolute beginners guide},
``Generating random numbers'', and \href{https://numpy.org/doc/stable/reference/random/index.html}%
{Random sampling}.

\textbf{In simplydata.} \code{sdata.make\_random\_numbers} makes normal or uniform numbers (and whole
numbers); \code{sdata.distribution("poisson", average=3).make\_random\_numbers(how\_many=10)} makes them
from any of 28 distributions. \code{numbers.make\_random\_picks} is \code{choice}. \code{with\_seed=}
is the seed.

\begin{lstlisting}[style=py]
heights = sdata.make_random_numbers(how_many=100, average=150, standard_deviation=8, with_seed=1)
dice = sdata.make_random_numbers(how_many=10, between=[1, 6], whole_numbers=True, with_seed=1)
arrivals = sdata.distribution("poisson", average=3).make_random_numbers(how_many=10, with_seed=1)
picked = heights.make_random_picks(how_many=5, with_seed=1)
print(picked)
\end{lstlisting}

\unsection{Saving and reading numbers}{from-numpy:files}

\textbf{The job.} Write numbers to a file and read them back (NumPy's \code{save}, \code{savetxt},
\code{loadtxt} and \code{genfromtxt}). In the guide: \href{https://numpy.org/doc/stable/user/basics.io.html}%
{I/O with NumPy}.

\textbf{In simplydata.} \code{numbers.save("file.csv")} writes CSV, Excel, Parquet, JSON or a sound
file, as the end of the name says. \code{sdata.number\_list("file.csv")} reads them back, and
\code{empty\_values=[-999]} is \code{genfromtxt}'s \code{missing\_values}. A netCDF file opens as a
block with \code{sdata.number\_block("file.nc")}.

\begin{lstlisting}[style=py]
numbers = sdata.number_list([4.1, 3.9, -999, 4.4])
numbers.save("readings.csv")
again = sdata.number_list("readings.csv", empty_values=[-999])
again.count_empty_numbers()
\end{lstlisting}

\unsection{Grids as matrices}{from-numpy:matrices}

\textbf{The job.} The determinant, the inverse and the eigenvalues of a square grid, and solving the
equations grid times x = b (NumPy's \code{np.linalg}). In the guides:
\href{https://numpy.org/doc/stable/user/absolute_beginners.html}{the absolute beginners guide},
``Working with mathematical formulas'', and
\href{https://numpy.org/doc/stable/reference/routines.linalg.html}{Linear algebra}.

\textbf{In simplydata.} \code{find\_determinant}, \code{take\_inverse} and \code{find\_eigenvalues};
\code{solve\_for(b)} is \code{solve}, and on a grid that is not square it gives the least-squares answer
(\code{lstsq}) and says so. \code{exactly=True} on each gives SymPy's exact answer, as fractions. To
solve equations, \code{solve\_for} is better than an inverse. \code{multiply\_as\_matrix\_by} is the
matrix product (\code{@}, \code{np.matmul}); \code{find\_rank}, \code{find\_size} and
\code{find\_condition\_number} are \code{matrix\_rank}, \code{norm} and \code{cond}.

\begin{lstlisting}[style=py]
grid = sdata.number_grid([[2, 1], [1, 3]])
grid.find_determinant()                  # 5
grid.find_eigenvalues()
x = grid.solve_for([5, 10])              # x = [1, 3]
print(x)
grid.solve_for([1, 1], exactly=True)     # [2/5, 1/5]
grid.multiply_as_matrix_by(sdata.number_grid([[0, 1], [1, 0]]))
grid.find_rank()                         # 2
\end{lstlisting}

\unsection{Fourier transforms}{from-numpy:fourier}

\textbf{The job.} Turn readings in time into frequencies, and back (NumPy's and SciPy's \code{fft},
\code{rfft}, \code{fftfreq} and \code{ifft}). In the guide:
\href{https://docs.scipy.org/doc/scipy/tutorial/fft.html}{Discrete Fourier Transforms (scipy.fft)}.

\textbf{In simplydata.} \code{numbers.take\_fourier\_transform(time\_step=)} gives a table with each
frequency in hertz, its size and its phase: \code{rfft} and \code{rfftfreq} in one word.
\code{sdata.find\_inverse\_fourier\_transform(table)} goes back. \code{take\_cosine\_transform} is
\code{dct}. \code{find\_spectrum} gives the power at each frequency.

\begin{lstlisting}[style=py]
wave = sdata.make_wave("sine", frequency=5, time_step=0.01, how_many=200)
frequencies = wave.take_fourier_transform(time_step=0.01)
back = sdata.find_inverse_fourier_transform(frequencies)
back.count()
\end{lstlisting}

\unsection{Areas, volumes and paths over time}{from-numpy:integration}

\textbf{The job.} The area under a curve, the volume under a surface, the running area under readings,
and how a quantity changes over time from a rate (SciPy's \code{quad}, \code{dblquad},
\code{cumulative\_trapezoid} and \code{solve\_ivp}). In the guide:
\href{https://docs.scipy.org/doc/scipy/tutorial/integrate.html}{Integration (scipy.integrate)}.

\textbf{In simplydata.} \code{sdata.find\_area\_under}, \code{sdata.find\_volume\_under},
\code{numbers.turn\_into\_running\_area} and \code{sdata.find\_path\_over\_time}. The formula is written
as text, with letters. \code{sdata.find\_path\_between\_ends} is \code{solve\_bvp}, for a path with a value
given at each end, and \code{sdata.find\_area\_points} gives the points and weights of Gaussian quadrature
(\code{roots\_legendre} and its kin).

\begin{lstlisting}[style=py]
sdata.find_area_under("exp(-x**2)", between=[0, "infinity"])
sdata.find_volume_under("x * y", between={"x": [0, 1], "y": [0, 2]})
flow = sdata.number_list([0.0, 1.5, 2.0, 1.0, 0.5])
flow.turn_into_running_area(time_step=60)
path = sdata.find_path_over_time("dN/dt = r * N", at_start={"N": 10}, between=[0, 10], known={"r": 0.3})
\end{lstlisting}

\unsection{Smallest, biggest and solutions}{from-numpy:optimization}

\textbf{The job.} Find where a formula is smallest or biggest, solve equations, fit a curve to data,
and solve least squares with limits (SciPy's \code{minimize}, \code{root}, \code{curve\_fit},
\code{linprog}, \code{nnls} and \code{lsq\_linear}). In the guide:
\href{https://docs.scipy.org/doc/scipy/tutorial/optimize.html}{Optimization (scipy.optimize)}.

\textbf{In simplydata.} \code{sdata.find\_where\_smallest} and \code{sdata.find\_where\_biggest} (with
\code{rules=} for linear programming), \code{sdata.solve}, \code{table.fit\_formula} for curves, and
\code{grid.solve\_for(b, between=[0, None])} for answers that may not be negative, as in a diet mix or
the unmixing of spectra. \code{sdata.find\_best\_pairing(costs=grid)} is \code{linear\_sum\_assignment}: it
pairs each row with one column so that the total cost is smallest.

\begin{lstlisting}[style=py]
sdata.find_where_smallest("x**2 - 4*x + 1", between=[0, 5])
sdata.find_where_smallest("2*x + 3*y", rules=["x + y >= 4", "x >= 0", "y >= 0"])
sdata.solve("x**3 - 2*x = 5", between=[0, 3])
# one row per nutrient (protein, fibre, fat), one column per food: grams of the nutrient in one portion
foods = sdata.number_grid([[20, 2, 5], [10, 15, 1], [2, 8, 30]])
amounts = foods.solve_for([60, 40, 50], between=[0, None])          # portions of each food, never below 0
grid = sdata.number_grid([[4, 1, 3], [2, 0, 5], [3, 2, 2]])         # cost of each worker (row) for each job
sdata.find_best_pairing(costs=grid)
\end{lstlisting}

\unsection{Values between the points}{from-numpy:interpolation}

\textbf{The job.} Find values between measured points, on a line, on a map or in a block (SciPy's
\code{interp1d}, \code{CubicSpline}, \code{griddata} and \code{RegularGridInterpolator}). In the guide:
\href{https://docs.scipy.org/doc/scipy/tutorial/interpolate.html}{Interpolation (scipy.interpolate)}.

\textbf{In simplydata.} \code{table.find\_values\_at} on a table's columns,
\code{table.make\_map\_of} for scattered points on a map, and \code{block.find\_values\_at} inside a
labelled block. \code{sdata.curve\_through\_points(x=..., y=...)} is a spline you can keep (SciPy's
\code{make\_interp\_spline} and \code{CubicSpline}): read it anywhere with \code{curve.find\_value\_at}, and
take its slope, area and zeros with \code{find\_slope\_at}, \code{find\_area\_between} and \code{find\_zeros}.

\begin{lstlisting}[style=py]
growth = sdata.get_example("growth")
growth.keep_rows_where("country", is_equal_to="Applia")
growth.find_values_at("population", at={"year": [2001.5, 2003.5]})
curve = sdata.curve_through_points(x=[0, 5, 10, 20], y=[8.1, 7.9, 6.2, 5.0])
curve.find_value_at(12)
curve.find_slope_at(12)
\end{lstlisting}

\unsection{Signals}{from-numpy:signals}

\textbf{The job.} Smooth readings, keep or remove frequencies with a filter, find peaks, and find the
spectrum (SciPy's \code{savgol\_filter}, \code{butter} with \code{filtfilt}, \code{find\_peaks} and
\code{welch}). In the guide: \href{https://docs.scipy.org/doc/scipy/tutorial/signal.html}{Signal
Processing (scipy.signal)}.

\textbf{In simplydata.} \code{smooth}, \code{keep\_each\_frequency\_that} and
\code{remove\_each\_frequency\_that}, \code{find\_each\_peak\_that} and \code{find\_spectrum}. Each
takes \code{time\_step=}, the time between readings, so frequencies are in hertz. The filter designs of
\code{scipy.signal} (Butterworth, Chebyshev, elliptic, Bessel and FIR filters) are values of
\code{sdata.frequency\_filter}; \code{system.move\_poles\_to} is \code{place\_poles}, and
\code{find\_phase\_locking\_to} measures how steady the phase is.

\begin{lstlisting}[style=py]
signal = sdata.make_wave("sine", frequency=2, time_step=0.01, how_many=300)
noise = sdata.make_random_numbers(how_many=300, average=0, standard_deviation=0.3, with_seed=1)
signal.add(noise)
signal.keep_each_frequency_that(is_below=5, time_step=0.01)
signal.find_each_peak_that(is_at_least=0.5)
\end{lstlisting}

\unsection{More linear algebra}{from-numpy:linalg}

\textbf{The job.} Break a matrix down into parts (SVD, QR, LU, Cholesky, Schur and more), the matrix
exponential, and the null space (SciPy's \code{scipy.linalg}). In the guide:
\href{https://docs.scipy.org/doc/scipy/tutorial/linalg.html}{Linear Algebra (scipy.linalg)}.

\textbf{In simplydata.} \code{grid.break\_down(test\_method="svd")} gives the parts as a dict of
grids and prints what each part is; the other methods are \code{"qr"}, \code{"lu"},
\code{"cholesky"}, \code{"schur"}, \code{"polar"}, \code{"ldl"}, \code{"hessenberg"} and \code{"qz"}.
\code{take\_matrix\_exp} is \code{expm} and \code{find\_null\_space} is \code{null\_space}. The other
matrix functions are one family word, \code{grid.take\_matrix\_function("log")} (\code{logm}, \code{sqrtm},
\code{cosm} ...); \code{check\_if\_grid\_is("symmetric")} checks a kind of grid, and
\code{sdata.solve\_matrix\_equation} solves the Sylvester, Lyapunov and Riccati equations.

\begin{lstlisting}[style=py]
grid = sdata.number_grid([[4, 2], [2, 3]])
parts = grid.break_down(test_method="cholesky")
print(parts["U"])
parts = grid.break_down(test_method="svd")
print(parts["singular values"])
grid.find_null_space()
grid.take_matrix_function("square root")                 # SciPy's sqrtm
grid.check_if_grid_is("symmetric")
\end{lstlisting}

\unsection{Points in space}{from-numpy:spatial}

\textbf{The job.} The nearest points, distances between points, and the fence round a set of points
(SciPy's \code{KDTree}, \code{distance.cdist} and \code{ConvexHull}). In the guide:
\href{https://docs.scipy.org/doc/scipy/tutorial/spatial.html}{Spatial data structures and algorithms
(scipy.spatial)}.

\textbf{In simplydata.} These are table words, because points are rows with an east and a north:
\code{table.find\_nearest\_rows\_to}, \code{table.find\_distances\_to\_rows\_of} and
\code{table.find\_fence\_around}. \code{table.find\_triangles\_between\_points} is \code{Delaunay} and
\code{table.find\_regions\_around\_points} is \code{Voronoi}, with the area of each region.

\begin{lstlisting}[style=py]
island = sdata.get_example("island")
island.find_nearest_rows_to({"east": 4, "north": 1}, how_many=3)
island.find_fence_around(x="east", y="north")
island = sdata.get_example("island")
island.find_triangles_between_points(x="east", y="north")
\end{lstlisting}

\unsection{Special functions}{from-numpy:special}

\textbf{The job.} The gamma and beta functions, the error function, Bessel functions, and the number of
ways to choose (SciPy's \code{scipy.special}). In the guide:
\href{https://docs.scipy.org/doc/scipy/tutorial/special.html}{Special Functions (scipy.special)}.

\textbf{In simplydata.} A formula can use \code{gamma}, \code{lgamma}, \code{digamma}, \code{beta},
\code{erf}, \code{erfc} and \code{besselJ}. \code{sdata.count\_ways\_to\_choose} is \code{comb} and
\code{sdata.count\_ways\_to\_order} is \code{perm}; \code{take\_log\_odds} and \code{turn\_into\_chances}
are \code{logit} and \code{expit}. Every other function of \code{scipy.special} is a name of
\code{numbers.take\_function} (see ``Maths on every number''), and the same names work in a formula with
\code{\_} for spaces, such as \code{bessel\_j(1, x)}. \code{sdata.find\_zeros\_of("bessel J")} gives the zeros
(\code{jn\_zeros} and its kin), and \code{sdata.polynomial("legendre", degree=3)} makes the polynomials of
\code{scipy.special} and \code{np.polynomial} as formulas.

\begin{lstlisting}[style=py]
formula = sdata.formula("gamma(x)")
formula.find_value_at(x=4.5)
sdata.count_ways_to_choose(3, out_of=10)                 # 120
chances = sdata.number_list([0.1, 0.5, 0.9])
chances.take_log_odds()
sdata.find_zeros_of("bessel J", how_many=3)
formula = sdata.polynomial("legendre", degree=3)
formula.find_value_at(x=0.5)
\end{lstlisting}

\unsection{Statistics}{from-numpy:statistics}

\textbf{The job.} Use a distribution (density, chances, percentiles, random numbers), fit one to data,
and run tests (SciPy's \code{scipy.stats}). In the guide:
\href{https://docs.scipy.org/doc/scipy/tutorial/stats.html}{Statistics (scipy.stats)}.

\textbf{In simplydata.} \code{sdata.distribution(...)} with \code{find\_density\_at} (\code{pdf}),
\code{find\_chance\_that\_value} (\code{cdf} and \code{sf}) and \code{find\_percentile} (\code{ppf});
\code{numbers.fit\_distribution} is \code{fit}. The tests are table words, such as
\code{table.compare\_averages\_of}; the rest of this manual explains them. Every distribution of
\code{scipy.stats} has a plain name, such as \code{sdata.distribution("skew normal", shape=4)}.

\begin{lstlisting}[style=py]
distribution = sdata.distribution("normal", average=50, standard_deviation=10)
distribution.find_chance_that_value(is_below=60)
distribution.find_percentile(95)
waits = sdata.make_random_numbers(how_many=200, average=10, standard_deviation=2, with_seed=1)
waits.fit_distribution("normal")
\end{lstlisting}

\unsection{Images and grids of cells}{from-numpy:ndimage}

\textbf{The job.} Smooth a grid, find patches of cells that pass a test, turn and resize a grid
(SciPy's \code{scipy.ndimage}: \code{gaussian\_filter}, \code{label}, \code{rotate} and \code{zoom}).
In the guide: \href{https://docs.scipy.org/doc/scipy/tutorial/ndimage.html}{Multidimensional image
processing (scipy.ndimage)}.

\textbf{In simplydata.} \code{grid.smooth}, \code{grid.find\_each\_patch\_that}, \code{grid.rotate} and
\code{grid.change\_size\_to}; on a block, \code{block.smooth(over=[...])}. \code{grid.find\_edges()} is
\code{sobel}; \code{grid.change\_patch\_shapes("grow")} is one family word for \code{binary\_dilation},
\code{binary\_erosion}, opening, closing and filling holes; \code{find\_distance\_to\_nearest} is
\code{distance\_transform\_edt}.

\begin{lstlisting}[style=py]
grid = sdata.number_grid([[0, 0, 1, 1], [0, 0, 1, 1], [5, 0, 0, 0], [5, 5, 0, 0]])
grid.find_each_patch_that(is_above=0)
grid.change_size_to(rows=8, columns=8)
grid.rotate(by_degrees=15)
grid = sdata.number_grid([[0, 0, 0, 0], [0, 1, 1, 0], [0, 1, 1, 0], [0, 0, 0, 0]])
grid.find_edges()
\end{lstlisting}

\unsection{Sound files and other files}{from-numpy:file-io}

\textbf{The job.} Read and write sound files, and other file types (SciPy's \code{scipy.io}, such as
\code{wavfile}). In the guide: \href{https://docs.scipy.org/doc/scipy/tutorial/io.html}{File IO
(scipy.io)}.

\textbf{In simplydata.} \code{numbers.save("sound.wav", time\_step=)} and
\code{sdata.number\_list("sound.wav")}. netCDF files open as a block, and
\code{sdata.show\_whats\_inside("file.nc")} lists what is in one.

\begin{lstlisting}[style=py]
tone = sdata.make_wave("sine", frequency=440, time_step=1 / 8000, how_many=8000)
tone.times(0.5)
tone.save("tone.wav", time_step=1 / 8000)
again = sdata.number_list("tone.wav")
again.count()
\end{lstlisting}

\unsection{The jobs left to NumPy and SciPy}{from-numpy:left-out}

simplydata does the jobs most people need in research and teaching, with one word each. Many jobs that
were once left to NumPy now have words: \code{np.diff} is \code{turn\_into\_changes}, \code{np.histogram} is
\code{count\_in\_bins}, \code{np.argsort} is \code{sort(show="positions")}, \code{np.meshgrid} is
\code{sdata.make\_coordinate\_grids}, \code{np.linalg.norm} is \code{find\_size}, the matrix product is
\code{multiply\_as\_matrix\_by}, and \code{np.polynomial} is \code{sdata.polynomial}. These jobs are left out
on purpose. Each one is a hand-over: \code{array = numbers.as\_numpy()} (or \code{grid.as\_numpy()},
\code{block.as\_numpy()}), then NumPy or SciPy code, then \code{sdata.number\_list(...)},
\code{sdata.number\_grid(...)} or \code{sdata.number\_block(...)} to come back.
\begin{itemize}
  \item \textbf{Sums over named indices} (\code{np.einsum}, \code{np.linalg.tensorinv}): einsum is its own
    notation; \code{multiply\_as\_matrix\_by} covers the usual uses.
  \item \textbf{Data types, copies and views, structured arrays and strings in arrays}: simplydata keeps
    numbers as decimals and makes copies, so these questions do not come up.
  \item \textbf{Sparse arrays} (\code{scipy.sparse}) and their eigenvalue and graph routines.
  \item \textbf{The intersection of half-spaces} (\code{scipy.spatial.HalfspaceIntersection}), expert
    geometry with no data job.
  \item \textbf{Steps inside a solver} (\code{qr\_update}, \code{expm\_frechet}, the BLAS and LAPACK
    functions): \code{sdata.help} names the word that runs them.
\end{itemize}

One more example: a matrix product, and back.

\begin{lstlisting}[style=py]
a = sdata.number_grid([[1, 2], [3, 4]])
b = sdata.number_grid([[0, 1], [1, 0]])
product = a.as_numpy() @ b.as_numpy()    # NumPy's matrix product
grid = sdata.number_grid(product)        # back to simplydata
grid.find_determinant()
\end{lstlisting}
